\section{The measure of the invariant tori (HEAD, HEAD-Q, G3)}\label{bp:p4}

% Local notation of Part 4 (prefixed to avoid clashes with other parts).
\newcommand{\pfh}{\mathbf{h}}
\newcommand{\pfD}{\mathfrak{B}}
\newcommand{\pfL}{L_{+}}
\newcommand{\pfg}{g}
\newcommand{\infd}{\operatorname{infDist}}

This part proves the three frozen theorems that concern measure: G3
(Theorem~\ref{bp:p0:g3}), HEAD (Theorem~\ref{bp:p0:head}) and HEAD-Q
(Theorem~\ref{bp:p0:headq}). The a-posteriori theorem G1 (Theorem~\ref{bp:p0:g1}) and local
uniqueness G2 (Theorem~\ref{bp:p0:g2}) are used as proved inputs, exactly as frozen; every
application lists the frozen hypotheses and the place where each is checked.

\paragraph{Route.} For a frequency $\omega=\nabla h(I_\omega)$ which is $(\gamma,\tau)$-Diophantine,
the unperturbed torus $K^0_\omega(\theta)=(\theta,I_\omega)$ is an approximate invariant torus of
$h+f$ with error $O(\varepsilon)$, and G1 turns it into an invariant torus $K_\omega$ at distance
$O(\varepsilon/\gamma^2)$ (Section~\ref{bp:p4:sec-apply}). A phase shift normalizes $K_\omega$
(Section~\ref{bp:p4:sec-normal}). For a second Diophantine frequency $\omega'$, the torus
$K_\omega$ shifted in the actions by $I_{\omega'}-I_\omega$ is an approximate invariant torus at
$\omega'$ whose error is $O(\varepsilon\gamma^{-2}\abs{\omega'-\omega})$; G1 and G2 then show that
$K_{\omega'}-K^0_{\omega'}$ and $K_\omega-K^0_\omega$ differ by $O(\varepsilon\gamma^{-4}\abs{\omega'-\omega})$
(Section~\ref{bp:p4:sec-lip}). Both re-applications need the hypotheses of G1 at tori that are
only close to $K^0$; Section~\ref{bp:p4:sec-stab} proves this stability. With
$\gamma\sim\varepsilon^{1/4}$, the map $(\theta,I)\mapsto K_{\nabla h(I)}(\theta)$ is a
Lipschitz-small perturbation of the identity on the set of good actions; a McShane extension makes
it a lipeomorphism of $\R^{2n}$, so the complement of the tori is the image of the complement of
the good actions, whose measure is $O(\gamma)$ by G3 and the boundary-layer hypothesis
(Section~\ref{bp:p4:sec-measure}).

\paragraph{Conventions of this part.} Norms are those of Part~\ref{bp:p0}: the max norm on
$\R^m$, $\C^m$ and on products, induced operator norms, and for matrices the row-sum norm
$\abs{A}=\max_i\sum_j\abs{A_{ij}}$. We write $\iota:\R^m\to\C^m$ for the real inclusion. The
Gram inverse called $B$ in Definition~\ref{bp:p0:stripdata} is written $G_K^{-1}$ here, with
$G_K=L^\top L$, since $B$ denotes the action domain. For a set $S\subset\R^m$, $\meas(S)$ is
Lebesgue outer measure (Lean \Lean{volume} applied to an arbitrary set); all measure inequalities
below are between outer measures and use only monotonicity and countable subadditivity, except
where measurability is stated.

%======================================================================
\subsection{G3: the measure of the non-Diophantine set}\label{bp:p4:sec-g3}

The set of frequencies in a bounded set that fail the $(\alpha,\tau)$-Diophantine condition is the
union over $k\ne0$ of slabs $\{\abs{k\cdot\omega}<\alpha\abs{k}_1^{-\tau}\}$. Each slab meets a
bounding cube in measure $O(\alpha\abs{k}_1^{-\tau-1})$, and the lattice sum
$\sum_k\abs{k}_1^{-\tau-1}$ converges exactly when $\tau>n-1$.
\source{Pos01 \S1d, p.~5: ``$m(R^\tau_{\alpha,k}\cap\Omega)=O(\alpha/\abs{k}^{\tau+1})$, and thus
$m(R^\tau_\alpha\cap\Omega)\le\sum_{k\ne0}m(R^\tau_{\alpha,k}\cap\Omega)=O(\alpha)$, provided that
$\tau>n-1$''. The two estimates are asserted there without proof; Lemmas~\ref{bp:p4:lattice-sum}
and~\ref{bp:p4:slab} supply them.}

\begin{lemma}[Counting lattice points on $\ell^1$ spheres]\label{bp:p4:lattice-count}
Let $n\ge1$ and $m\ge1$ be integers. Then $\#\{k\in\Z^n:\abs{k}_1=m\}\le2(2m+1)^{n-1}\le2\cdot3^{n-1}m^{n-1}$.
\end{lemma}

\begin{proof}
\proofstep{Projection to the first $n-1$ coordinates}\label{step:p4:count-projection}
Let $P(k)=(k_1,\dots,k_{n-1})$. If $\abs{k}_1=m$ then $\abs{k_i}\le m$, so $P(k)$ lies in
$\{-m,\dots,m\}^{n-1}$, a set with $(2m+1)^{n-1}$ elements (one element if $n=1$). The fibre of $P$
over $(k_1,\dots,k_{n-1})$ is contained in $\{\pm(m-\sum_{i<n}\abs{k_i})\}$, since
$\abs{k_n}=m-\sum_{i<n}\abs{k_i}$; it has at most two elements. Hence the count is at most
$2(2m+1)^{n-1}$, and $2m+1\le3m$ for $m\ge1$.
\end{proof}

\begin{lemma}[Lattice sum]\label{bp:p4:lattice-sum}
Let $n\ge1$ and $\tau>n-1$. Then
\begin{equation*}
\sum_{k\in\Z^n\setminus\{0\}}\abs{k}_1^{-\tau-1}\le S(n,\tau):=2\cdot3^{n-1}\,\frac{\tau+2-n}{\tau+1-n}\,.
\end{equation*}
\end{lemma}

\begin{proof}
\proofstep{Summation by shells}\label{step:p4:latsum-shells}
Group the terms by $m=\abs{k}_1\ge1$ and apply Lemma~\ref{bp:p4:lattice-count}: the sum is at
most $2\cdot3^{n-1}\sum_{m\ge1}m^{-p}$ with $p=\tau+2-n>1$.

\proofstep{Integral comparison}\label{step:p4:latsum-integral}
Since $x\mapsto x^{-p}$ decreases, $m^{-p}\le\int_{m-1}^mx^{-p}\,dx$ for $m\ge2$, so
$\sum_{m\ge1}m^{-p}\le1+\int_1^\infty x^{-p}\,dx=1+\frac1{p-1}=\frac{p}{p-1}$. With
$p=\tau+2-n$ this is the claim.
\end{proof}

\begin{lemma}[Slabs in a cube]\label{bp:p4:slab}
Let $n\ge1$, $R>0$, $a>0$ and $k\in\Z^n\setminus\{0\}$. Then
\begin{equation*}
\meas\{\omega\in[-R,R]^n:\abs{k\cdot\omega}<a\}\le\frac{2n(2R)^{n-1}a}{\abs{k}_1}\,.
\end{equation*}
\end{lemma}

\begin{proof}
\proofstep{A dominant coordinate}\label{step:p4:slab-coordinate}
Choose $i$ with $\abs{k_i}=\max_j\abs{k_j}$. Then $\abs{k_i}\ge1$ and $\abs{k}_1\le n\abs{k_i}$.

\proofstep{Fubini in the coordinate $\omega_i$}\label{step:p4:slab-fubini}
The slab is open, hence measurable. Fix the other coordinates
$\omega'=(\omega_j)_{j\ne i}\in[-R,R]^{n-1}$ and put $b=\sum_{j\ne i}k_j\omega_j$. The section
$\{\omega_i\in[-R,R]:\abs{k_i\omega_i+b}<a\}$ is contained in the interval
$(-b/k_i-a/\abs{k_i},-b/k_i+a/\abs{k_i})$, of length $2a/\abs{k_i}$. By Fubini's theorem (Lebesgue
measure on $\R^n$ is the product of the one-dimensional measures), the slab has measure at most
$(2R)^{n-1}\cdot2a/\abs{k_i}\le2n(2R)^{n-1}a/\abs{k}_1$.
\end{proof}

\begin{proof}[Proof of Theorem~\ref{bp:p0:g3}]
\proofstep{The case $n=0$}\label{step:p4:g3-n0}
If $n=0$, the only $k\in\Z^0$ is $k=0$, so every $\omega$ is $(\alpha,\tau)$-Diophantine, the set
is empty, and any $C$ works.

\proofstep{Bounding cube and constant}\label{step:p4:g3-cube}
Let $n\ge1$. Since $\Omega$ is bounded, choose $R>0$ with $\Omega\subset[-R,R]^n$ and put
$C=2n(2R)^{n-1}S(n,\tau)$, with $S(n,\tau)$ from Lemma~\ref{bp:p4:lattice-sum}. The constant
depends only on $n$, $\tau$ and $R$.

\proofstep{Union of slabs}\label{step:p4:g3-union}
Let $\alpha>0$. By Definition~\ref{bp:p0:diophantine}, $\omega\notin\Dio_{\alpha,\tau}$ means that
$\abs{k\cdot\omega}<\alpha\abs{k}_1^{-\tau}$ for some $k\ne0$ (here $\abs{k}_1\ge1$, so the real
power is the usual one). Hence
\begin{equation*}
\{\omega\in\Omega:\omega\notin\Dio_{\alpha,\tau}\}\subset\bigcup_{k\ne0}\{\omega\in[-R,R]^n:\abs{k\cdot\omega}<\alpha\abs{k}_1^{-\tau}\}\,.
\end{equation*}

\proofstep{Summation}\label{step:p4:g3-sum}
By monotonicity and countable subadditivity of the outer measure, Lemma~\ref{bp:p4:slab} with
$a=\alpha\abs{k}_1^{-\tau}$, and Lemma~\ref{bp:p4:lattice-sum},
\begin{equation*}
\meas\{\omega\in\Omega:\omega\notin\Dio_{\alpha,\tau}\}\le\sum_{k\ne0}2n(2R)^{n-1}\alpha\abs{k}_1^{-\tau-1}\le C\alpha\,.
\end{equation*}
\end{proof}

\begin{remark}\label{bp:p4:g3-remark}
The proof gives $C\ge0$, and the constant depends on $\Omega$ only through a bounding cube. This is
Lemma~\ref{bp:p1:dmeas}. Below, G3 is applied once, to the bounded set $\nabla h(B)$ of
Lemma~\ref{bp:p4:freq-map}.
\end{remark}

%======================================================================
\subsection{Uniform data from the hypotheses of HEAD}\label{bp:p4:sec-data}

Throughout Sections~\ref{bp:p4:sec-data}--\ref{bp:p4:sec-measure} we fix the data of HEAD:
$n\ge2$, $\tau>n-1$, $s_0,r_0>0$, an open bounded set $B\subset\R^n$, a function $h$ holomorphic on
$B_{r_0}$ and real at real points of $B_{r_0}$, a witness $L\in\R$ of the nondegeneracy hypothesis
and a witness $C_B\in\R$ of the boundary-layer hypothesis. We assume $B\ne\emptyset$; the empty
case is trivial (Step~\ref{step:p4:head-empty}). All constants named in this part are computed from
these data alone, never from $f$ or $\varepsilon$; a constant ``depending on the data'' means
depending only on $n,\tau,s_0,r_0,B,h,L,C_B$. The perturbation $f$ and the number $\varepsilon$
are as in HEAD: $f$ is holomorphic on $D(s_0,r_0)$, $1$-periodic in $\theta$, real at real points,
and $\abs f\le\varepsilon$ on $D(s_0,r_0)$. We put
\begin{equation*}
\pfh(\theta,I)=h(I)+f(\theta,I)\,,\qquad (\theta,I)\in D(s_0,r_0)\,,
\end{equation*}
so that the frozen real Hamiltonian is $H(z)=\operatorname{Re}\pfh(\iota z)$ at every real point
$z$ with $\iota z\in D(s_0,r_0)$ (Definition~\ref{bp:p0:nearlyintegrable}).

\begin{definition}[Shrunken domains and basic constants]\label{bp:p4:data}
Put $\pfL=\max(L,1)$, $C_B^+=\max(C_B,0)$, and
\begin{equation*}
s=\min(s_0,1)\,,\qquad \delta_*=\tfrac14\min(s,r_0)\,,\qquad \mu=\tfrac12\delta_*\,,\qquad
\pfD=T_{s/2}\times B_{r_0/2}\,,\qquad \pfD'=T_{3s/4}\times B_{3r_0/4}\,.
\end{equation*}
Let $c(m,j)$ be the constants of Lemma~\ref{bp:p1:cauchy-open} and
$\bar c=\max\{1,c(n,j),c(2n,j):j=1,2,3\}$. For $I_*\in\R^n$ the \emph{unperturbed torus at $I_*$}
is $K^0_{I_*}(\theta)=(\theta,\iota I_*)$, $\theta\in\C^n$.
\end{definition}

Note $\delta_*\le\frac14$, $\mu\le\frac18$, and $\pfD\subset\pfD'\subset D(s_0,r_0)$ because
$s\le s_0$.

\begin{lemma}[$h$ is bounded on a shrunken neighbourhood]\label{bp:p4:h-bounded}
$M_h:=\sup\{\abs{h(I)}:I\in B_{3r_0/4}\}$ is finite and depends only on the data.
\end{lemma}

\begin{proof}
\proofstep{A compact set between the neighbourhoods}\label{step:p4:hbdd-compact}
Let $Q=\{I\in\C^n:\infd(I,\iota B)\le\frac34r_0\}$. It is closed, and bounded because $B$ is
bounded, hence compact. It contains $B_{3r_0/4}$. If $I\in Q$ then some $J\in B$ has
$\abs{I-\iota J}<r_0$, since $\frac34r_0<r_0$; so $Q\subset B_{r_0}$.

\proofstep{Continuity}\label{step:p4:hbdd-continuity}
$h$ is holomorphic, hence continuous, on $B_{r_0}\supset Q$, so $\abs h$ is bounded on the
compact set $Q$ and a fortiori on $B_{3r_0/4}$.
\end{proof}

\begin{lemma}[Derivative bounds]\label{bp:p4:derivative-bounds}
Put $H_j=\bar cM_h\delta_*^{-j}$ and $F_j=\bar c\delta_*^{-j}$ for $j=1,2,3$, and
$M=\bar c(M_h+1)\delta_*^{-3}$.
\begin{enumerate}[label=(\alph*)]
\item For $z\in\pfD$, the closed polydisc of radius $\delta_*$ around $z$ lies in $\pfD'$; for
  $I\in B_{r_0/2}$, the closed polydisc of radius $\delta_*$ around $I$ lies in $B_{3r_0/4}$.
\item $\abs{D^jh(I)}\le H_j$ for $I\in B_{r_0/2}$ and $j=1,2,3$.
\item $\abs{D^jf(z)}\le F_j\varepsilon$ for $z\in\pfD$ and $j=1,2,3$; in particular
  $\abs{X_f}\le F_1\varepsilon$ and $\abs{DX_f}\le F_2\varepsilon$ on $\pfD$.
\item If $\varepsilon\le1$, then on $\pfD$:
  $\abs{D\pfh},\abs{X_\pfh},\abs{DX_\pfh},\abs{D^2X_\pfh}\le M$.
\end{enumerate}
\end{lemma}

\begin{proof}
\proofstep{Polydiscs}\label{step:p4:db-polydisc}
Let $z=(\theta,I)\in\pfD$ and $\abs{w-z}\le\delta_*$, $w=(\theta',I')$. Then
$\abs{\operatorname{Im}\theta'_i}\le\abs{\operatorname{Im}\theta_i}+\delta_*<\frac s2+\frac s4$, and
for $J\in B$ with $\abs{I-\iota J}<\frac{r_0}2$ we get $\abs{I'-\iota J}<\frac{r_0}2+\delta_*\le\frac34r_0$.
The second claim of (a) is the action half of this computation.

\proofstep{Cauchy estimates}\label{step:p4:db-cauchy}
On $B_{3r_0/4}$, $\abs h\le M_h$ (Lemma~\ref{bp:p4:h-bounded}); on $\pfD'\subset D(s_0,r_0)$,
$\abs f\le\varepsilon$ and $\abs\pfh\le M_h+\varepsilon$. Lemma~\ref{bp:p1:cauchy-open} with the
polydiscs of (a) gives (b), (c) for $D^jf$, and $\abs{D^j\pfh}\le\bar c(M_h+\varepsilon)\delta_*^{-j}$ on
$\pfD$; since $\delta_*\le1$, $\delta_*^{-j}\le\delta_*^{-3}$.

\proofstep{From $D\pfh$ to $X_\pfh$}\label{step:p4:db-vectorfield}
For a holomorphic function $\phi$ on an open subset of $\C^n\times\C^n$, the frozen definition
(\Lean{hamiltonianVectorField\ensuremath{\mathbb{C}}}) reads $X_\phi=J\circ D\phi$, where
$J:\mathcal L(\C^{2n},\C)\to\C^{2n}$, $J(\ell)=\bigl((\ell(0,e_i))_i,(-\ell(e_i,0))_i\bigr)$, is
linear with $\abs{J(\ell)}\le\norm{\ell}$ because $\abs{(0,e_i)}=\abs{(e_i,0)}=1$. Since $J$ is a
continuous linear map, $DX_\phi=J\circ D^2\phi$ and $D^2X_\phi=J\circ D^3\phi$, whence
$\abs{X_\phi}\le\abs{D\phi}$, $\abs{DX_\phi}\le\abs{D^2\phi}$, $\abs{D^2X_\phi}\le\abs{D^3\phi}$.
With $\phi=f$ this completes (c); with $\phi=\pfh$ and $\varepsilon\le1$ it gives (d), since
$\abs{D\pfh}\le\bar c(M_h+1)\delta_*^{-1}\le M$. (In Lean, the norm of
\Lean{iteratedFDeriv} $2$ of $X_\phi$ equals that of the second derivative of $X_\phi$, and
composition with $J$ is \Lean{ContinuousLinearMap.compContinuousMultilinearMap}.)
\end{proof}

\begin{lemma}[The frequency map]\label{bp:p4:freq-map}
Let $U=\{I\in\R^n:\iota I\in B_{r_0}\}$ and $\pfg=\Lean{frequencyMap}\,h:\R^n\to\R^n$.
\begin{enumerate}[label=(\roman*)]
\item $U$ is open and contains $\bar B$.
\item For $I\in U$ the complex partial derivatives $\partial_ih(\iota I)$ are real and
  $\pfg(I)=\nabla h(\iota I):=(\partial_ih(\iota I))_{i}$.
\item $\pfg$ is differentiable on $U$, and $D\pfg(I)$ is the real symmetric matrix
  $\nabla^2h(\iota I)=(\partial_j\partial_ih(\iota I))_{ij}$.
\item For $\iota I\in B_{r_0/2}$: $\abs{\pfg(I)}\le H_1$ and $\abs{\nabla^2h(\iota I)}\le H_2$.
\item $\pfg$ is $\Lambda$-Lipschitz on $B$, with $\Lambda=\max(H_2,2H_1/\delta_*)$, and
  $\pfg(B)$ is bounded by $W:=H_1$.
\end{enumerate}
\end{lemma}

\begin{proof}
\proofstep{The real domain}\label{step:p4:fm-domain}
$U=\iota^{-1}(B_{r_0})$ is open as the preimage of an open set. If $J\in\bar B$, some $I\in B$ has
$\abs{J-I}<r_0$, so $\iota J\in B_{r_0}$.

\proofstep{Guarded fact: the frozen frequency map is the gradient}\label{step:p4:fm-gradient}
Let $\varphi(J)=h(\iota J)$ for $J\in U$. As the composition of the $\C$-differentiable $h$ with
the $\R$-linear $\iota$, $\varphi$ is $\R$-differentiable on $U$ with
$D\varphi(J)v=Dh(\iota J)(\iota v)$. By hypothesis $\varphi$ is real-valued on the open set $U$,
so $D\varphi(J)v=\lim_{t\to0}(\varphi(J+tv)-\varphi(J))/t$ is real; in particular
$\partial_ih(\iota J)=D\varphi(J)e_i\in\R$. The function $J\mapsto\operatorname{Re}h(\iota J)$ in the
definition of $\pfg$ agrees with $\varphi$ on the neighbourhood $U$ of $I$, so its Fréchet
derivative at $I$ is $D\varphi(I)$, and $\pfg_i(I)=D\varphi(I)e_i=\partial_ih(\iota I)$. This
establishes the guarded fact of Remark~\ref{bp:p0:totality} on $U$.

\proofstep{The Hessian}\label{step:p4:fm-hessian}
Each $\partial_ih$ is holomorphic on $B_{r_0}$ and real on $\iota U$ by (ii). The argument of
Step~\ref{step:p4:fm-gradient} applied to $\partial_ih$ shows that $\pfg_i$ is differentiable on
$U$ with $D\pfg_i(I)v=\sum_j\partial_j\partial_ih(\iota I)v_j$, all entries real; the matrix is
symmetric because mixed complex partial derivatives of holomorphic functions commute.

\proofstep{Bounds}\label{step:p4:fm-bounds}
$\abs{\pfg(I)}=\max_i\abs{Dh(\iota I)e_i}\le\abs{Dh(\iota I)}\le H_1$ and, row by row,
$\sum_j\abs{\partial_j\partial_ih}=\sup_{\abs v\le1}\abs{D^2h[e_i,v]}\le\abs{D^2h}\le H_2$, by
Lemma~\ref{bp:p4:derivative-bounds}(b).

\proofstep{Lipschitz bound on $B$}\label{step:p4:fm-lipschitz}
Let $I,J\in B$. If $\abs{I-J}<\delta_*$, every point $x$ of the segment $[I,J]$ has
$\abs{x-I}<\delta_*<\frac{r_0}2$, so $\iota x\in B_{r_0/2}$; the mean value inequality and (iv)
give $\abs{\pfg(I)-\pfg(J)}\le H_2\abs{I-J}$. If $\abs{I-J}\ge\delta_*$, then
$\abs{\pfg(I)-\pfg(J)}\le2H_1\le(2H_1/\delta_*)\abs{I-J}$. Boundedness of $\pfg(B)$ is (iv).
\end{proof}

\begin{lemma}[Kolmogorov nondegeneracy at a point]\label{bp:p4:hessian-inverse}
For every $I\in B$ and $v\in\R^n$, $\abs v\le\pfL\abs{\nabla^2h(\iota I)v}$. Consequently the
matrix $\nabla^2h(\iota I)$ is invertible, and its inverse is a real matrix with
$\abs{\nabla^2h(\iota I)^{-1}}\le\pfL$ in the row-sum norm, also as a complex matrix.
\end{lemma}

\begin{proof}
\proofstep{Difference quotients}\label{step:p4:hess-quotient}
Since $L\le\pfL$, the hypothesis gives $\abs{I-J}\le\pfL\abs{\pfg(I)-\pfg(J)}$ for $I,J\in\bar B$.
As $B$ is open, $I+tv\in B$ for small $t>0$, so
$\abs v\le\pfL\abs{(\pfg(I+tv)-\pfg(I))/t}$. Letting $t\to0^+$ and using
Lemma~\ref{bp:p4:freq-map}(iii) gives $\abs v\le\pfL\abs{\nabla^2h(\iota I)v}$.

\proofstep{Inverse and its norm}\label{step:p4:hess-norm}
A real square matrix $A$ with $\abs v\le\pfL\abs{Av}$ is injective, hence invertible, and
$\abs{A^{-1}w}\le\pfL\abs w$ for real $w$. The row-sum norm of a real matrix $P$ is
$\max_i\sum_j\abs{P_{ij}}=\max_i\abs{(Pw^{(i)})_i}$ with $w^{(i)}_j=\operatorname{sign}P_{ij}$,
$\abs{w^{(i)}}\le1$; hence $\abs{A^{-1}}\le\pfL$, and this is also the complex operator norm.
\end{proof}

\begin{lemma}[The inverse frequency map]\label{bp:p4:inverse-freq}
Let $\Omega=\pfg(B)$. The map $\pfg$ is injective on $\bar B$, and its inverse
$\pfg(\bar B)\to\bar B$ is $\pfL$-Lipschitz. For $\omega\in\Omega$ let $I_\omega$ be the unique
$I\in B$ with $\pfg(I)=\omega$. Then $I_\omega\in B$, $\pfg(I_\omega)=\omega$,
$\abs{I_\omega-I_{\omega'}}\le\pfL\abs{\omega-\omega'}$, and $\abs\omega\le W$. We write
$K^0_\omega=K^0_{I_\omega}$.
\end{lemma}

\begin{proof}
\proofstep{Inverse Lipschitz bound}\label{step:p4:invfreq-proof}
The inequality $\abs{I-J}\le\pfL\abs{\pfg(I)-\pfg(J)}$ on $\bar B$
(Step~\ref{step:p4:hess-quotient}) gives injectivity and the Lipschitz bound of the inverse;
$\abs\omega\le W$ is Lemma~\ref{bp:p4:freq-map}(v).
\end{proof}

\begin{lemma}[The Hamiltonian hypothesis H1]\label{bp:p4:h1}
Let $\varepsilon\le1$. Then $\pfD$ and $\pfh$ satisfy the hypotheses on $(\mathfrak B,h)$ of G1
and G2 (the block ``H1'' of Theorems~\ref{bp:p0:g1} and~\ref{bp:p0:g2}) with the constant $M$ of
Lemma~\ref{bp:p4:derivative-bounds}.
\end{lemma}

\begin{proof}
\proofstep{Openness and invariance}\label{step:p4:h1-open}
$T_{s/2}$ is open (finitely many strict inequalities between continuous functions) and invariant
under $\theta\mapsto\theta+k$, $k\in\Z^n$, because $\operatorname{Im}(\theta+k)=\operatorname{Im}\theta$;
$B_{r_0/2}$ is a union of open balls. So $\pfD$ is open and $\Z^n$-invariant in $\theta$.

\proofstep{Holomorphy, periodicity, reality}\label{step:p4:h1-holo}
$(\theta,I)\mapsto h(I)$ is holomorphic on $\C^n\times B_{r_0}$ and $f$ on
$D(s_0,r_0)\supset\pfD$, so $\pfh$ is holomorphic on $\pfD$. For $z\in\pfD$ and $k\in\Z^n$,
$z\in D(s_0,r_0)$, so $f(\theta+k,I)=f(\theta,I)$ by the periodicity hypothesis of HEAD, and $h(I)$
does not depend on $\theta$. If $\iota(\theta,I)\in\pfD$, then $\iota I\in B_{r_0}$ and
$\iota(\theta,I)\in D(s_0,r_0)$, so $h(\iota I)$ and $f(\iota\theta,\iota I)$ are real by the
reality hypotheses of HEAD.

\proofstep{Derivative bounds}\label{step:p4:h1-bounds}
Lemma~\ref{bp:p4:derivative-bounds}(d).
\end{proof}

%======================================================================
\subsection{The invariant torus at each Diophantine frequency (APPLY)}\label{bp:p4:sec-apply}

For a Diophantine $\omega\in\pfg(B)$, the flat torus $K^0_\omega(\theta)=(\theta,I_\omega)$ is
invariant for $h$ alone; for $h+f$ its invariance error is the vector field of $f$, of size
$O(\varepsilon)$. Its tangent frame is the identity and its torsion is the Hessian of $h$ plus an
$O(\varepsilon)$ average, so the hypotheses of G1 hold with constants depending on the data only.

\begin{lemma}[Geometry of the unperturbed torus]\label{bp:p4:unperturbed}
Let $I_*\in B$, $\omega\in\R^n$, and $K^0=K^0_{I_*}$. Assume $\varepsilon\le1$.
\begin{enumerate}[label=(\alph*)]
\item $K^0$ is holomorphic on $\C^n$, primary on every strip, and real on reals;
  $DK^0(\theta)=L^0:=\left(\begin{smallmatrix}\mathrm{Id}\\0\end{smallmatrix}\right)$ and $\abs{L^0}=1$.
\item The Gram matrix is $G_{K^0}=\mathrm{Id}$; its determinant is a unit and
  $\abs{G_{K^0}^{-1}}=1$.
\item The normal frame is constant: $N^0(\theta)e_j=(0,e_j)$, i.e.
  $N^0=\left(\begin{smallmatrix}0\\\mathrm{Id}\end{smallmatrix}\right)$, and $DN^0(\theta)[\omega]=0$.
\item For $\theta\in T_{s/2}$ the torsion is
  $T^0(\theta)=\partial_I^2\pfh(\theta,\iota I_*)=\nabla^2h(\iota I_*)+\partial_I^2f(\theta,\iota I_*)$,
  whatever $\omega$ is. It is continuous, so $\avg{T^0}$ is the integral of a continuous function,
  and $\avg{T^0}=\nabla^2h(\iota I_*)+\avg{\partial_I^2f(\cdot,\iota I_*)}$.
\item If $\pfL F_2\varepsilon\le\frac12$, then $\avg{T^0}$ is invertible and
  $\abs{\avg{T^0}^{-1}}\le2\pfL$.
\item If $\omega=\pfg(I_*)$, the invariance error at $\omega$ is $E^0(\theta)=X_f(\theta,\iota I_*)$
  for $\theta\in T_{s/2}$, and $\abs{E^0}\le F_1\varepsilon$ there.
\item For $\theta\in T_{s/4}$, the open ball of radius $\delta_*$ around $K^0(\theta)$ lies in
  $\pfD$.
\end{enumerate}
\end{lemma}

\begin{proof}
\proofstep{Frame and Gram matrix}\label{step:p4:unp-frame}
$K^0$ is affine, so $DK^0(\theta)v=(v,0)$ for all $\theta$; $K^0(\theta+k)=K^0(\theta)+(k,0)$ for
$k\in\Z^n$, and $K^0(\iota x)=\iota(x,I_*)$ for real $x$. Each row of $L^0$ has one entry $1$ or
none, so $\abs{L^0}=1$. With the bilinear dot product of the frozen \Lean{gramMatrix},
$G_{ij}=\Lean{dot\ensuremath{\mathbb{C}}}((e_i,0),(e_j,0))=\delta_{ij}$. The determinant is $1$,
and the Lean inverse of the identity matrix is the identity.

\proofstep{Normal frame}\label{step:p4:unp-normal}
By the frozen \Lean{normalFrame}, $N^0(\theta)e_j=\sum_k(G^{-1})_{kj}\Omega_0(e_k,0)=\Omega_0(e_j,0)=(0,e_j)$,
since $\Omega_0(a,b)=(-b,a)$. This does not depend on $\theta$, so its derivative in any direction
vanishes.

\proofstep{Torsion}\label{step:p4:unp-torsion}
Let $z=(\theta,\iota I_*)\in\pfD$. By Step~\ref{step:p4:db-vectorfield},
$DX_\pfh(z)(0,e_j)=(a,b)$ with $a_m=D^2\pfh(z)[(0,e_j),(0,e_m)]$ and
$b_m=-D^2\pfh(z)[(0,e_j),(e_m,0)]$. The frozen torsion is
$T^0_{ij}=\Lean{dot\ensuremath{\mathbb{C}}}\bigl((0,e_i),\Omega_0(DX_\pfh(z)(0,e_j)-0)\bigr)=\Lean{dot\ensuremath{\mathbb{C}}}((0,e_i),(-b,a))=a_i=\partial_{I_j}\partial_{I_i}\pfh(z)$.
The derivative term vanishes by Step~\ref{step:p4:unp-normal}, so $\omega$ does not enter. The
function $\theta\mapsto T^0(\theta)$ is holomorphic on $T_{s/2}$, hence continuous on the real unit
cube, and the frozen \Lean{torsionAverage} is a genuine integral of it; integrating the constant
$\nabla^2h(\iota I_*)$ over the unit cube returns it.

\proofstep{Twist}\label{step:p4:unp-twist}
Row by row, $\sum_j\abs{\partial_{I_j}\partial_{I_i}f(z)}\le\abs{D^2f(z)}\le F_2\varepsilon$ by
Lemma~\ref{bp:p4:derivative-bounds}(c); integrating over the unit cube (measure $1$) bounds the
row-sum norm of $\Delta=\avg{\partial^2_If(\cdot,\iota I_*)}$ by $F_2\varepsilon$. By
Lemma~\ref{bp:p4:hessian-inverse}, $A=\nabla^2h(\iota I_*)$ is invertible with $\abs{A^{-1}}\le\pfL$,
and $\abs{A^{-1}}\abs\Delta\le\frac12$. Lemma~\ref{bp:p1:neumann} gives (e).

\proofstep{Invariance error}\label{step:p4:unp-error}
Write $\pfh=\tilde h+f$ with $\tilde h(\theta,I)=h(I)$. Then $X_{\tilde h}(\theta,I)=(\nabla h(I),0)$,
and $\nabla h(\iota I_*)=\pfg(I_*)=\omega$ by Lemma~\ref{bp:p4:freq-map}(ii). Since
$DK^0(\theta)\omega=(\omega,0)$, the frozen \Lean{invarianceError} is
$X_\pfh(z)-(\omega,0)=X_f(z)$, bounded by $F_1\varepsilon$ by
Lemma~\ref{bp:p4:derivative-bounds}(c).

\proofstep{Neighbourhood}\label{step:p4:unp-ball}
If $\abs{(\theta',I')-(\theta,\iota I_*)}<\delta_*$ with $\theta\in T_{s/4}$, then
$\abs{\operatorname{Im}\theta'_i}<\frac s4+\delta_*\le\frac s2$ and
$\abs{I'-\iota I_*}<\delta_*<\frac{r_0}2$ with $I_*\in B$.
\end{proof}

\begin{definition}[Parameters of the applications of G1 and G2]\label{bp:p4:parameters}
Put
\begin{equation*}
\sigma_K=\sigma_B=2+\mu\,,\qquad\sigma_T=4\pfL+\mu\,,\qquad\bar\gamma=1\,.
\end{equation*}
Theorem~\ref{bp:p0:g1} is applied with two tuples $(n,\tau,\rho,\rho_\infty,M,\sigma_K,\sigma_B,\sigma_T,\mu,\bar\gamma)$,
namely $(\rho,\rho_\infty)=(\frac s4,\frac s8)$, with constants $C_1^{(1)},C_2^{(1)}$, and
$(\rho,\rho_\infty)=(\frac{3s}{32},\frac s{16})$, with constants $C_1^{(2)},C_2^{(2)}$.
Theorem~\ref{bp:p0:g2} is applied with the tuple $(n,\tau,\rho,M,\sigma_K,\sigma_B,\sigma_T,\mu)$,
$\rho=\frac s{16}$, with constant $c_G$. The side conditions of the frozen statements hold:
$n\ge2$, $\tau\ge n-1$, $0<\rho_\infty<\rho$, and $M,\sigma_K,\sigma_B,\sigma_T,\mu,\bar\gamma>0$.
The derived constants are
\begin{gather*}
A_0=\frac{2C_2^{(1)}F_1}{(s/4)^{2\tau}}\,,\quad
\lambda_*=\frac{\delta_*}{4\pfL}\,,\quad
A_1=\pfL(H_3A_0+F_2)+\frac{32nA_0}{s}\,,\quad
A_2=\frac{(\sigma_K+1)C_2^{(2)}A_1}{(3s/32)^{2\tau}}\,,\\
A_3=\max\Bigl(A_2,\frac{2A_0}{\lambda_*}\Bigr)\,,\quad
A_4=\frac{16nA_0}{s}\,,\quad
c_T=304n^3(M+W+1)\,,\quad
\eta_*(\delta)=\min\Bigl(\mu,\frac{\delta^2}{16c_T\pfL}\Bigr)\,,
\end{gather*}
and $c_*>0$ is the minimum of the eight numbers
\begin{gather*}
1\,,\quad\frac{\delta_*^2}{2\pfL\bar c}\,,\quad\frac{(s/4)^{4\tau}}{2C_1^{(1)}F_1}\,,\quad
\frac{\eta_*(s/32)}{A_0}\,,\quad\frac{(3s/32)^{4\tau}}{2C_1^{(2)}A_1\lambda_*}\,,\\
\frac{c_G(s/16)^{2\tau}}{2A_0+C_2^{(2)}A_1\lambda_*(3s/32)^{-2\tau}}\,,\quad
\frac1{4(A_3\Lambda+A_4)}\,,\quad\frac1{2A_0}\,.
\end{gather*}
All of these depend only on the data.
\end{definition}

The second entry is $\delta_*^2/(2\pfL\bar c)=1/(2\pfL F_2)$. The \emph{standing assumption} of
Sections~\ref{bp:p4:sec-apply}--\ref{bp:p4:sec-real} is
\begin{equation}\label{bp:p4:SA}
0<\gamma\le1\qquad\text{and}\qquad0\le\varepsilon\le c_*\gamma^4\,.
\end{equation}
We write $\eta_0=A_0\varepsilon\gamma^{-2}$ and $\Omega_\gamma=\pfg(B)\cap\Dio_{\gamma,\tau}$.

\begin{lemma}[Smallness consequences]\label{bp:p4:smallness}
Under~\eqref{bp:p4:SA}:
\begin{enumerate}[label=(s\arabic*)]
\item $\varepsilon\le1$ and $\pfL F_2\varepsilon\le\frac12$;
\item $C_1^{(1)}F_1\varepsilon/(\gamma^4(s/4)^{4\tau})\le\frac12$;
\item $\eta_0\le\eta_*(\frac s{32})\le\eta_*(\frac s{16})$, $\eta_0\le\mu$, and $16n\eta_0/s\le\frac12$;
\item $C_1^{(2)}A_1\lambda_*\varepsilon/(\gamma^4(3s/32)^{4\tau})\le\frac12$;
\item $\bigl(2A_0+C_2^{(2)}A_1\lambda_*(3s/32)^{-2\tau}\bigr)\varepsilon\gamma^{-2}\le c_G\gamma^2(s/16)^{2\tau}$;
\item $(A_3\Lambda+A_4)\varepsilon\gamma^{-4}\le\frac14$;
\item $\eta_0\le\frac12\gamma^2\le\frac12\min(\gamma,1)$.
\end{enumerate}
\end{lemma}

\begin{proof}
\proofstep{Reduction to $c_*$}\label{step:p4:small-proof}
Since $\gamma\le1$, $\varepsilon\le c_*\gamma^4\le c_*\gamma^2\le c_*$. Each item is then the
corresponding entry of the minimum defining $c_*$: (s1) from the first two entries, (s2), (s4),
(s6) directly, (s5) and (s7) after multiplying $\varepsilon\gamma^{-2}\le c_*\gamma^2$ by the
relevant constant, and (s3) from $\eta_0\le c_*A_0\gamma^2\le\eta_*(s/32)$, the monotonicity of
$\eta_*$, $\eta_*\le\mu$, and $16n\eta_*(s/32)/s\le ns/(1024c_T)\le\frac12$ because
$c_T\ge n$ and $s\le1$.
\end{proof}

\begin{theorem}[APPLY: a torus for every Diophantine frequency]\label{bp:p4:apply}
Assume~\eqref{bp:p4:SA} and let $\omega\in\Omega_\gamma$. There is
$K^\infty_\omega:\C^n\to\C^n\times\C^n$ which is holomorphic on $T_{s/8}$, primary on $T_{s/8}$,
real on reals, invariant with frequency $\omega$ on $T_{s/8}$ ($E=0$ there), with
$B_{\mu/2}(K^\infty_\omega(\theta))\subset\pfD$, $\abs{DK^\infty_\omega(\theta)}\le\sigma_K-\mu/2$ and
$\abs{K^\infty_\omega(\theta)-K^0_\omega(\theta)}\le\frac12\eta_0$ for $\theta\in T_{s/8}$.
\end{theorem}

\begin{proof}
\proofstep{Application of G1 at the unperturbed torus}\label{step:p4:apply-g1}
Apply Theorem~\ref{bp:p0:g1} with the first tuple of Definition~\ref{bp:p4:parameters},
$\mathfrak B=\pfD$, $h=\pfh$, $K=K^0_\omega$, the frequency $\omega$, the given $\gamma$, and
$e=F_1\varepsilon$. The frozen hypotheses are checked as follows.
\begin{itemize}
\item H1 ($\mathfrak B$ open and $\Z^n$-invariant; $\pfh$ holomorphic, periodic, real; the four
  derivative bounds by $M$): Lemma~\ref{bp:p4:h1}, applicable by (s1).
\item H2 ($K$ holomorphic on $T_{s/4}$, primary, real on reals; $B_\mu(K(\theta))\subset\mathfrak B$;
  $\abs{DK}\le\sigma_K-\mu=2$): Lemma~\ref{bp:p4:unperturbed}(a), and (g) with $\mu<\delta_*$.
\item H3 (Gram determinant a unit, $\abs{G^{-1}}\le\sigma_B-\mu=2$): Lemma~\ref{bp:p4:unperturbed}(b).
\item H4 ($\avg T$ invertible, $\abs{\avg T^{-1}}\le\sigma_T-\mu=4\pfL$): Lemma~\ref{bp:p4:unperturbed}(e),
  applicable because $I_\omega\in B$ and by (s1).
\item H5 ($0<\gamma\le\bar\gamma=1$, $\omega\in\Dio_{\gamma,\tau}$): \eqref{bp:p4:SA} and
  $\omega\in\Omega_\gamma$.
\item Error ($\abs{E}\le e$ on $T_{s/4}$ and $C_1^{(1)}e/(\gamma^4(s/4)^{4\tau})<1$):
  Lemma~\ref{bp:p4:unperturbed}(f) with $I_*=I_\omega$, $\pfg(I_\omega)=\omega$, and (s2).
\end{itemize}
G1 yields $K^\infty_\omega$ with all the stated properties on $T_{s/8}$, and
$\abs{K^\infty_\omega-K^0_\omega}\le C_2^{(1)}F_1\varepsilon/(\gamma^2(s/4)^{2\tau})=\frac12A_0\varepsilon\gamma^{-2}$.
\end{proof}

\begin{remark}[Why no distance to the boundary of $\pfg(B)$]\label{bp:p4:omega-gamma}
Pöschel's set $\Omega_\alpha$ (Pos01, p.~6) keeps the frequencies at distance $\alpha$ from the
boundary of the frequency domain. Here G1 needs no such condition: $\pfD$ contains the
$\delta_*$-neighbourhood of every $K^0_{I_*}$ with $I_*\in B$ (Lemma~\ref{bp:p4:unperturbed}(g)),
because the analyticity domain $B_{r_0}$ is a uniform neighbourhood of $B$. The boundary layer
enters only in the measure step, in action space, where the hypothesis of HEAD lives
(Section~\ref{bp:p4:sec-measure}); the frequency domain $\pfg(B)$ need not have a regular boundary.
\end{remark}

%======================================================================
\subsection{Stability of the hypotheses of G1 (STAB)}\label{bp:p4:sec-stab}

The Lipschitz argument re-applies G1 at a torus produced by G1 and uses such a torus as the
reference torus of G2. Both tori are only known to be $C^0$-close to an unperturbed torus
$K^0_{I_*}$ on a strip. Cauchy estimates turn $C^0$-closeness into closeness of the first two
derivatives on a thinner strip, and the Neumann series then transports the bounds of
Lemma~\ref{bp:p4:unperturbed} on the Gram inverse and the twist. We always compare with an
unperturbed torus, never with another perturbed one, so the bounds do not deteriorate from one
application to the next.

\begin{lemma}[Matrix norms]\label{bp:p4:matrix-norms}
\begin{enumerate}[label=(\roman*)]
\item For $A\in\C^{p\times q}$, the row-sum norm $\abs A$ is the operator norm of $A$ from
  $(\C^q,\abs\cdot_\infty)$ to $(\C^p,\abs\cdot_\infty)$. In particular it is submultiplicative,
  and for a linear map $\ell:\C^q\to\C^p$ (such as \Lean{fderiv} of $K$ at a point) the Lean
  operator norm $\norm\ell$ is the row-sum norm of its matrix.
\item $\abs{A^\top}\le p\abs A$ for $A\in\C^{p\times q}$.
\item $\abs{\Omega_0}=\abs{L^0}=\abs{(L^0)^\top}=\abs{N^0}=\abs{(N^0)^\top}=1$, with $L^0,N^0$
  as in Lemma~\ref{bp:p4:unperturbed}.
\item In the bilinear notation of Definition~\ref{bp:p0:stripdata},
  $\Lean{dot\ensuremath{\mathbb{C}}}(x,\Omega_0y)=x^\top\Omega_0y$, $\Omega_0^\top\Omega_0=\mathrm{Id}$,
  $\Omega_0^2=-\mathrm{Id}$, and for real $x,y$,
  $\Lean{symplecticForm}(x,y)=(\iota x)^\top\Omega_0(\iota y)$.
\item If $A:[0,1]^n\to\C^{p\times q}$ is continuous, then $\abs{\int A}\le\sup\abs A$.
\end{enumerate}
\end{lemma}

\begin{proof}
\proofstep{Row sums}\label{step:p4:mn-rowsum}
Items (i) and (ii) and $\abs{\Omega_0}=1$ are Lemmas~\ref{bp:p1:operator-norms}
and~\ref{bp:p1:matrix-norms}; for (ii) directly: a row sum of $A^\top$ is a column sum of $A$, at
most $p$ times the largest entry, which is at most $\abs A$. For (iii), each of these matrices has in each
row at most one nonzero entry, of modulus $1$, and some row has one. For (iv), compute with
$\Omega_0(a,b)=(-b,a)$: $x^\top\Omega_0y=\sum_i(x_{2,i}y_{1,i}-x_{1,i}y_{2,i})$, which is the frozen
formula of \Lean{symplecticForm} for real arguments. Item (v) is
Lemma~\ref{bp:p1:averages-matrix}: integrate the row sums over the unit cube, of measure $1$.
\end{proof}

\begin{lemma}[Guarded facts for the frame and the torsion]\label{bp:p4:torsion-regular}
Let $\rho>0$, let $K$ be holomorphic on $T_\rho$ with $K(T_\rho)\subset\pfD$ and
$\det G_K(\theta)\ne0$ for $\theta\in T_\rho$, and let $\omega\in\R^n$. Write $L=DK$ and
$\partial_\omega F=DF\,\omega$.
\begin{enumerate}[label=(\alph*)]
\item The Lean matrix inverse $G_K(\theta)^{-1}$ is the inverse of $G_K(\theta)$, and
  $\theta\mapsto G_K(\theta)^{-1}$ is holomorphic on $T_\rho$ with
  $\partial_\omega(G_K^{-1})=-G_K^{-1}(\partial_\omega G_K)G_K^{-1}$.
\item $N=\Omega_0LG_K^{-1}$ is holomorphic on $T_\rho$, and the derivative term of the frozen
  \Lean{torsion} is the directional derivative $\partial_\omega N$.
\item On $T_\rho$, $T=N^\top\Omega_0X_N$ with $X_N=DX_\pfh(K)N-\partial_\omega N$, and $T$ is
  holomorphic, so $\avg T$ is the integral of a continuous function on the real unit cube.
\item If $K$ is primary on $T_\rho$, then $u(\theta)=K_1(\theta)-\theta$, $v(\theta)=K_2(\theta)$,
  $L$, $G_K$, $N$ and $T$ are $\Z^n$-periodic on $T_\rho$.
\end{enumerate}
\end{lemma}

\begin{proof}
\proofstep{Inverse}\label{step:p4:tr-inverse}
$L=DK$ is holomorphic on the open set $T_\rho$, and the entries of $G_K$ and its determinant are
polynomials in the entries of $L$. The Lean inverse is $\operatorname{Ring.inverse}(\det G)$ times
the adjugate; where $\det G\ne0$ this is $(\det G)^{-1}\operatorname{adj}G=G^{-1}$. Holomorphy and
the derivative formula are Lemma~\ref{bp:p1:holomorphic-inverse}.

\proofstep{Frame and torsion}\label{step:p4:tr-torsion}
By (a), $N$ is holomorphic, so its Fréchet derivative (the \Lean{fderiv} in the frozen
\Lean{torsion}) is its true derivative. $X_\pfh$ is holomorphic on the open set $\pfD\ni K(\theta)$,
so $DX_\pfh(K(\theta))$ is its true derivative there. Writing the frozen sums as matrix products
with Lemma~\ref{bp:p4:matrix-norms}(iv) gives $T=N^\top\Omega_0X_N$, a holomorphic function.

\proofstep{Periodicity}\label{step:p4:tr-periodic}
If $K(\theta+k)=K(\theta)+(k,0)$ on $T_\rho$, then $u(\theta+k)=K_1(\theta)+k-\theta-k=u(\theta)$
and $v(\theta+k)=v(\theta)$; differentiating the identity on the open set $T_\rho$ gives
$L(\theta+k)=L(\theta)$, and $G_K,N,T$ are built from $L$, its derivative and $K$ through the
$\Z^n$-periodic $X_\pfh$ (Lemma~\ref{bp:p4:h1}).
\end{proof}

\begin{theorem}[STAB]\label{bp:p4:stab}
Assume $\varepsilon\le1$ and $\pfL F_2\varepsilon\le\frac12$. Let $I_*\in B$, $\omega\in\R^n$
with $\abs\omega\le W$, $0<\delta<\rho\le\frac s4$, and let $K$ be holomorphic on $T_\rho$ with
\begin{equation*}
\abs{K(\theta)-K^0_{I_*}(\theta)}\le\eta\quad(\theta\in T_\rho)\,,\qquad\eta\le\eta_*(\delta)\,.
\end{equation*}
Then for $\theta\in T_{\rho-\delta}$:
\begin{enumerate}[label=(\alph*)]
\item $\abs{DK(\theta)-L^0}\le n\eta/\delta$ and $\abs{DK(\theta)}\le2=\sigma_K-\mu$;
\item $B_\mu(K(\theta))\subset\pfD$;
\item $\det G_K(\theta)\ne0$ and $\abs{G_K(\theta)^{-1}}\le2=\sigma_B-\mu$;
\item the torsion of $K$ at $\omega$, computed on $T_{\rho-\delta}$, satisfies
  $\abs{T(\theta)-T^0(\theta)}\le c_T\eta\delta^{-2}$, where $T^0$ is the torsion of $K^0_{I_*}$;
  $\avg T$ is invertible and $\abs{\avg T^{-1}}\le4\pfL=\sigma_T-\mu$.
\end{enumerate}
If in addition $K$ is primary on $T_\rho$ and real on reals, then $K$ satisfies the hypotheses
H2--H4 of G1 on $T_{\rho-\delta}$, relative to $(\pfD,\pfh)$, with the parameters
$\sigma_K,\sigma_B,\sigma_T,\mu$ of Definition~\ref{bp:p4:parameters}.
\end{theorem}

\begin{proof}
\proofstep{Smallness bookkeeping}\label{step:p4:stab-small}
Put $\Delta=K-K^0_{I_*}$, $a=n\eta/\delta$ and $b=4n^2W\eta\delta^{-2}$. Since $\delta<\frac s4\le1$
and $\eta\le\delta^2/(16c_T\pfL)$ with $c_T\ge304n^3$, we have $a\le n\eta\delta^{-2}\le\frac1{10n}$
and $c_T\eta\delta^{-2}\le\frac1{16}$.

\proofstep{Cauchy estimates}\label{step:p4:stab-cauchy}
By Lemma~\ref{bp:p1:cauchy-strip} on $T_\rho$ with loss $\delta$, $\abs{D\Delta}\le a$ on
$T_{\rho-\delta}$. With loss $\delta/2$ twice, $\abs{\partial_\ell\Delta}\le2\eta/\delta$ on
$T_{\rho-\delta/2}$ and $\abs{\partial_m\partial_\ell\Delta}\le4\eta/\delta^2$ on $T_{\rho-\delta}$. Since
$D^2K^0_{I_*}=0$, the matrix $\partial_\omega L=\bigl(\sum_m\omega_m\partial_m\partial_\ell\Delta\bigr)_\ell$
has row sums at most $n\cdot n\cdot W\cdot4\eta\delta^{-2}$, so $\abs{\partial_\omega L}\le b$ on
$T_{\rho-\delta}$. In particular $\abs L\le\abs{L^0}+a\le2$, which is (a).

\proofstep{Neighbourhood}\label{step:p4:stab-ball}
For $\theta\in T_{\rho-\delta}\subset T_{s/4}$, $\abs{K(\theta)-K^0_{I_*}(\theta)}\le\eta\le\mu$, so
$B_\mu(K(\theta))\subset B_{2\mu}(K^0_{I_*}(\theta))=B_{\delta_*}(K^0_{I_*}(\theta))\subset\pfD$ by
Lemma~\ref{bp:p4:unperturbed}(g). This is (b), and it gives $K(T_{\rho-\delta})\subset\pfD$.

\proofstep{Gram matrix}\label{step:p4:stab-gram}
With $L=L^0+D\Delta$ and Lemma~\ref{bp:p4:matrix-norms}(ii),(iii),
$\abs{G_K-\mathrm{Id}}\le\abs{(L^0)^\top}a+2na\abs{L^0}+2na^2\le5na\le\frac12$. By
Lemma~\ref{bp:p1:neumann} with $A=\mathrm{Id}$, $G_K$ is invertible (so $\det G_K\ne0$),
$\abs{G_K^{-1}}\le2$ and $\abs{G_K^{-1}-\mathrm{Id}}\le10na$. This is (c), and
Lemma~\ref{bp:p4:torsion-regular} applies on $T_{\rho-\delta}$.

\proofstep{Normal frame}\label{step:p4:stab-frame}
$N-N^0=\Omega_0\bigl((L-L^0)G_K^{-1}+L^0(G_K^{-1}-\mathrm{Id})\bigr)$, so
$\abs{N-N^0}\le2a+10na\le12na$ and $\abs N\le2$. By Lemma~\ref{bp:p4:torsion-regular}(a),
$\abs{\partial_\omega G_K}\le\abs{(\partial_\omega L)^\top}\abs L+\abs{L^\top}\abs{\partial_\omega L}\le8nb$,
$\abs{\partial_\omega(G_K^{-1})}\le32nb$, and
$\abs{\partial_\omega N}\le\abs{\partial_\omega L}\abs{G_K^{-1}}+\abs L\abs{\partial_\omega(G_K^{-1})}\le66nb$.

\proofstep{The vector $X_N$}\label{step:p4:stab-xn}
The segment from $K^0_{I_*}(\theta)$ to $K(\theta)$ lies in the convex ball of radius $\delta_*$
around $K^0_{I_*}(\theta)$, which is contained in $\pfD$, where $\abs{D^2X_\pfh}\le M$
(Lemma~\ref{bp:p4:h1}); by Lemma~\ref{bp:p1:increments-open} applied to $DX_\pfh$ on this ball,
$\abs{DX_\pfh(K)-DX_\pfh(K^0_{I_*})}\le M\eta$. With $X_N^0=DX_\pfh(K^0_{I_*})N^0$ (the torsion
vector of $K^0_{I_*}$, Lemma~\ref{bp:p4:unperturbed}(c)),
\begin{equation*}
\abs{X_N-X_N^0}\le M\eta\abs N+M\abs{N-N^0}+\abs{\partial_\omega N}\le(2M+12n^2M+264n^3W)\eta\delta^{-2}\le c_T\eta\delta^{-2}\,,
\end{equation*}
using $\delta\le1$. Hence $\abs{X_N}\le M+1$.

\proofstep{Torsion}\label{step:p4:stab-torsion}
$T-T^0=(N-N^0)^\top\Omega_0X_N+(N^0)^\top\Omega_0(X_N-X_N^0)$, so by
Lemma~\ref{bp:p4:matrix-norms},
$\abs{T-T^0}\le2n\cdot12na(M+1)+(2M+12n^2M+264n^3W)\eta\delta^{-2}\le c_T\eta\delta^{-2}$.

\proofstep{Twist}\label{step:p4:stab-twist}
By Lemma~\ref{bp:p4:torsion-regular}(c) and Lemma~\ref{bp:p4:unperturbed}(d), both averages are
integrals of continuous functions, and Lemma~\ref{bp:p4:matrix-norms}(v) gives
$\abs{\avg T-\avg{T^0}}\le c_T\eta\delta^{-2}$. By Lemma~\ref{bp:p4:unperturbed}(e),
$\abs{\avg{T^0}^{-1}}\le2\pfL$, and $2\pfL c_T\eta\delta^{-2}\le\frac18$. Lemma~\ref{bp:p1:neumann}
gives (d).

\proofstep{The hypotheses H2--H4}\label{step:p4:stab-h2h4}
H2: holomorphy, primarity and reality restrict from $T_\rho$ to $T_{\rho-\delta}$; the ball and
the derivative bound are (b) and (a). H3 is (c), a nonzero complex determinant being a unit. H4
is (d). The integration domain of $\avg T$ is the real unit cube, contained in $T_{\rho-\delta}$.
\end{proof}

%======================================================================
\subsection{Phase normalization}\label{bp:p4:sec-normal}

An invariant torus is determined by its image only up to a phase: $\theta\mapsto K(\theta+a)$
parameterizes the same set. G2 identifies two nearby invariant tori only up to such a phase. We
remove the freedom by requiring that the angle displacement have zero average; the Lipschitz
estimate of Section~\ref{bp:p4:sec-lip} compares normalized tori.
\source{CCdlL13 \S5.1.1, normalization (5.3) and Rem.~25--27, p.~31; Cor.~30 and its proof, p.~33
(``It is important that we choose $K$ normalized'').}

\begin{definition}[Angle displacement and phase average]\label{bp:p4:phase-average}
Let $K$ be holomorphic on $T_\rho$, primary on $T_\rho$ and real on reals. Its \emph{angle
displacement} is $u_K(\theta)=K_1(\theta)-\theta$ and its \emph{phase average} is
$\operatorname{ph}(K)=\int_{[0,1]^n}u_K(\iota x)\,dx\in\R^n$; it is real because $K_1(\iota x)$ is
real for real $x$, and finite because $u_K$ is continuous. $K$ is \emph{normalized} if
$\operatorname{ph}(K)=0$. For $a\in\R^n$ put $K^{[a]}(\theta)=K(\theta+\iota a)$.
\end{definition}

\begin{lemma}[Averages of periodic functions are translation invariant]\label{bp:p4:periodic-average}
If $\phi:\R^n\to\C^m$ is continuous and $\Z^n$-periodic, then
$\int_{[0,1]^n}\phi(x+a)\,dx=\int_{[0,1]^n}\phi(x)\,dx$ for every $a\in\R^n$.
\end{lemma}

\begin{proof}
\proofstep{One variable at a time}\label{step:p4:pa-proof}
By Fubini it suffices to treat $n=1$. For a continuous $1$-periodic $\psi$, the function
$a\mapsto\int_0^1\psi(x+a)\,dx=\int_a^{a+1}\psi$ has derivative $\psi(a+1)-\psi(a)=0$.
(Mathlib: \Lean{Function.Periodic.intervalIntegral\_add\_eq}.)
\end{proof}

\begin{lemma}[Phase shifts]\label{bp:p4:shift}
Let $K$ be holomorphic, primary on $T_\rho$, real on reals, $a\in\R^n$, and $\omega\in\R^n$.
\begin{enumerate}[label=(\alph*)]
\item $K^{[a]}$ is holomorphic on $T_\rho$, primary on $T_\rho$ and real on reals, and
  $DK^{[a]}(\theta)=DK(\theta+\iota a)$ for $\theta\in T_\rho$.
\item The invariance error at $\omega$ satisfies $E_{K^{[a]}}(\theta)=E_K(\theta+\iota a)$; if
  $K$ is invariant on $T_\rho$, so is $K^{[a]}$.
\item Pointwise bounds on $T_\rho$ (on $\abs{DK}$, on balls around $K(\theta)$) transfer from $K$ to
  $K^{[a]}$.
\item $\operatorname{ph}(K^{[a]})=\operatorname{ph}(K)+a$. In particular a torus has exactly one
  normalized phase shift, $K^{[-\operatorname{ph}(K)]}$, and if $K$ and $K^{[a]}$ are both
  normalized then $a=0$.
\item For $I_*\in\R^n$: $\abs{\operatorname{ph}(K)}\le\sup_{x\in\R^n}\abs{K(\iota x)-K^0_{I_*}(\iota x)}$
  and $\abs{K^{[a]}(\theta)-K^0_{I_*}(\theta)}\le\abs{K(\theta+\iota a)-K^0_{I_*}(\theta+\iota a)}+\abs a$.
\end{enumerate}
\end{lemma}

\begin{proof}
\proofstep{Translation}\label{step:p4:shift-translation}
$T_\rho$ is invariant under real translations, so $K^{[a]}$ is holomorphic there with the stated
derivative (chain rule). $K^{[a]}(\theta+k)=K(\theta+\iota a+k)=K^{[a]}(\theta)+(k,0)$, and
$K^{[a]}(\iota x)=K(\iota(x+a))$ is real. Then
$E_{K^{[a]}}(\theta)=X_\pfh(K(\theta+\iota a))-DK(\theta+\iota a)\omega=E_K(\theta+\iota a)$, which
gives (b), and (c) is the same substitution.

\proofstep{Phase average}\label{step:p4:shift-phase}
$u_{K^{[a]}}(\iota x)=K_1(\iota(x+a))-\iota x=u_K(\iota(x+a))+\iota a$, and $u_K\circ\iota$ is continuous
and $\Z^n$-periodic (Lemma~\ref{bp:p4:torsion-regular}(d)); Lemma~\ref{bp:p4:periodic-average}
gives (d). For (e), $u_{K^0_{I_*}}=0$, so $u_K(\iota x)$ is the angle component of
$K(\iota x)-K^0_{I_*}(\iota x)$, and the average over a set of measure $1$ is bounded by the sup;
the second bound is the triangle inequality with
$\abs{K^0_{I_*}(\theta+\iota a)-K^0_{I_*}(\theta)}=\abs a$.
\end{proof}

\begin{definition}[The normalized torus $K_\omega$]\label{bp:p4:normalized-torus}
Under~\eqref{bp:p4:SA}, for $\omega\in\Omega_\gamma$ let $K^\infty_\omega$ be a torus given by
Theorem~\ref{bp:p4:apply} (fixed by choice) and put
$K_\omega=(K^\infty_\omega)^{[-\operatorname{ph}(K^\infty_\omega)]}$. Write
$K_\omega=K^0_\omega+W_\omega$ and $W_\omega=(u_\omega,v_\omega)$, so that
$K_\omega(\theta)=(\theta+u_\omega(\theta),\iota I_\omega+v_\omega(\theta))$.
\end{definition}

\begin{lemma}[Properties of $K_\omega$]\label{bp:p4:normalize}
Under~\eqref{bp:p4:SA}, for $\omega\in\Omega_\gamma$: $K_\omega$ is holomorphic on $T_{s/8}$,
primary on $T_{s/8}$, real on reals, invariant with frequency $\omega$ on $T_{s/8}$, normalized
($\avg{u_\omega}=0$), $B_{\mu/2}(K_\omega(\theta))\subset\pfD$,
$\abs{DK_\omega(\theta)}\le\sigma_K-\mu/2$ and $\abs{W_\omega(\theta)}\le\eta_0$ for
$\theta\in T_{s/8}$.
\end{lemma}

\begin{proof}
\proofstep{Shift of the G1 output}\label{step:p4:normalize-proof}
Everything except the last bound is Theorem~\ref{bp:p4:apply} with Lemma~\ref{bp:p4:shift}(a)--(d).
By Lemma~\ref{bp:p4:shift}(e) and Theorem~\ref{bp:p4:apply},
$\abs{\operatorname{ph}(K^\infty_\omega)}\le\frac12\eta_0$ and
$\abs{W_\omega}\le\frac12\eta_0+\frac12\eta_0$.
\end{proof}

\begin{corollary}[Uniqueness of the normalized torus]\label{bp:p4:normalized-unique}
Assume~\eqref{bp:p4:SA} and $\omega\in\Omega_\gamma$. Let $K'$ be holomorphic on $T_{s/8}$,
primary on $T_{s/8}$, real on reals, invariant with frequency $\omega$ on $T_{s/8}$, normalized,
and $\abs{K'-K^0_\omega}\le\eta_0$ on $T_{s/8}$. Then $K'=K_\omega$ on $T_{s/8}$. In particular
$K_\omega$ does not depend on the choice made in Definition~\ref{bp:p4:normalized-torus}.
\end{corollary}

\begin{proof}
\proofstep{Application of G2}\label{step:p4:nu-g2}
Apply Theorem~\ref{bp:p0:g2} with the tuple of Definition~\ref{bp:p4:parameters} ($\rho=\frac s{16}$),
$\mathfrak B=\pfD$, $h=\pfh$, $K_1=K_\omega$, $K_2=K'$, the frequency $\omega$ and $\gamma$.
\begin{itemize}
\item H1: Lemma~\ref{bp:p4:h1} and (s1).
\item $K_1$ satisfies H2--H4 on $T_{s/16}$: Theorem~\ref{bp:p4:stab} with $I_*=I_\omega\in B$,
  $\abs\omega\le W$, $\rho=\frac s8$, $\delta=\frac s{16}$, $\eta=\eta_0\le\eta_*(\frac s{16})$ by
  (s3); $K_\omega$ is primary and real on reals (Lemma~\ref{bp:p4:normalize}).
\item $0<\gamma$, $\omega\in\Dio_{\gamma,\tau}$, and $E_1=0$ on $T_{s/16}$: Lemma~\ref{bp:p4:normalize}.
\item $K_2$ is holomorphic on $T_{s/16}$, primary, real on reals, $E_2=0$ there; and
  $K_2(T_{s/16})\subset\pfD$ because $\abs{K'(\theta)-K^0_\omega(\theta)}\le\eta_0<\delta_*$ and
  Lemma~\ref{bp:p4:unperturbed}(g).
\item Closeness: $\abs{K_1-K_2}\le2\eta_0=2A_0\varepsilon\gamma^{-2}\le c_G\gamma^2(s/16)^{2\tau}$
  by (s5).
\end{itemize}
G2 gives $\varphi\in\R^n$ with $K'(\theta)=K_\omega(\theta+\iota\varphi)$ for $\theta\in T_{s/32}$.

\proofstep{The phase vanishes}\label{step:p4:nu-phase}
On real points, $K'=K_\omega^{[\varphi]}$, and the phase average only involves real points. By
Lemma~\ref{bp:p4:shift}(d), $0=\operatorname{ph}(K')=\operatorname{ph}(K_\omega)+\varphi=\varphi$. So
$K'=K_\omega$ on $T_{s/32}$, and by the identity theorem on the connected open set $T_{s/8}$,
$K'=K_\omega$ on $T_{s/8}$.
\end{proof}

%======================================================================
\subsection{Lipschitz dependence on the frequency (LIP)}\label{bp:p4:sec-lip}

We compare $K_{\omega'}$ with $K_\omega$ by re-applying G1 at $\omega'$ to the torus $K_\omega$
shifted in the actions by $I_{\omega'}-I_\omega$. Its invariance error at $\omega'$ has three
parts: the change of $X_h$, which is a mixed second difference of $\nabla h$ and so carries the
factor $\abs{v_\omega}\abs{I_{\omega'}-I_\omega}$; the change of $X_f$, of size
$\varepsilon\abs{I_{\omega'}-I_\omega}$; and the transport term $-Du_\omega\,(\omega'-\omega)$. The
pure $\abs{\omega'-\omega}^2$ term cancels because the shift is chosen exactly so that
$\nabla h(I_{\omega'})-\nabla h(I_\omega)=\omega'-\omega$. G1 and G2 then compare the result with
$K_{\omega'}$, and normalization fixes the phase.
\source{PLAN, ``Route verdict'', item Lipschitz; template CCdlL13 Cor.~29--30, pp.~32--33 (ledger
D-LIP-1).}
\sourcegap{CCdlL13 Cor.~30 (p.~33) re-applies Theorem~20 at the solution for every pair
$\omega,\omega'$ (``we can apply the argument again to all $\omega,\omega'\in\tilde D_s$'') and states
that the non-degeneracy conditions ``depend only on $K_0,\mu_0$ and therefore they are uniform''.
That remark covers the first application only; the re-application at a computed solution needs the
stability of the non-degeneracy conditions, which is Theorem~\ref{bp:p4:stab} here. Moreover the
Lipschitz constant in Cor.~30 is not small, while the measure argument needs a small one; the
action shift below is our supplement.}

\begin{lemma}[Mixed second differences]\label{bp:p4:mixed-difference}
Let $U\subset\C^n$ be open, $F:U\to\C^n$ holomorphic, and $a,v,w\in\C^n$ with
$a+sv+tw\in U$ for all $s,t\in[0,1]$. If $\abs{D^2F}\le Q$ at these points, then
\begin{equation*}
\abs{F(a+v+w)-F(a+v)-F(a+w)+F(a)}\le Q\abs v\abs w\,.
\end{equation*}
\end{lemma}

\begin{proof}
\proofstep{Two mean value inequalities}\label{step:p4:mixed-proof}
Let $\psi(s)=F(a+sv+w)-F(a+sv)$ for $s\in[0,1]$. The left side is $\abs{\psi(1)-\psi(0)}$, and
$\psi'(s)=DF(a+sv+w)v-DF(a+sv)v$. The segment $t\mapsto a+sv+tw$ lies in $U$, so by the mean value
inequality applied to $DF$, $\abs{\psi'(s)}\le Q\abs w\abs v$. A second mean value inequality on
$[0,1]$ gives the claim.
\end{proof}

\begin{lemma}[The error of the shifted torus]\label{bp:p4:lip-error}
Assume~\eqref{bp:p4:SA}, let $\omega,\omega'\in\Omega_\gamma$ with $\abs{\omega'-\omega}\le\lambda_*\gamma^2$,
and put $\Delta\omega=\omega'-\omega$, $\Delta I=I_{\omega'}-I_\omega$ and
$K'=K_\omega+(0,\iota\Delta I)$. Then $K'$ is holomorphic on $T_{s/8}$, primary on $T_{s/8}$, real
on reals, $K'-K^0_{\omega'}=W_\omega$ (so $\abs{K'-K^0_{\omega'}}\le\eta_0$ on $T_{s/8}$), and the
invariance error $E'$ of $K'$ at $\omega'$ satisfies
\begin{equation*}
\abs{E'(\theta)}\le A_1\varepsilon\gamma^{-2}\abs{\Delta\omega}\qquad(\theta\in T_{3s/32})\,.
\end{equation*}
\end{lemma}

\begin{proof}
\proofstep{Elementary properties}\label{step:p4:le-basic}
Adding the real constant $(0,\iota\Delta I)$ preserves holomorphy, primarity and reality, and
$K^0_{\omega'}=K^0_\omega+(0,\iota\Delta I)$. By Lemma~\ref{bp:p4:inverse-freq},
$\abs{\Delta I}\le\pfL\abs{\Delta\omega}\le\pfL\lambda_*\gamma^2\le\frac14\delta_*$, and
$\abs{W_\omega}\le\eta_0\le\frac12\delta_*$ by (s3). Hence $\abs{W_\omega}+\abs{\Delta I}\le\frac34\delta_*$.

\proofstep{Splitting the error}\label{step:p4:le-split}
For $\theta\in T_{s/8}$, $X_\pfh(K_\omega(\theta))=DK_\omega(\theta)\omega$ (Lemma~\ref{bp:p4:normalize}),
and $DK_\omega\Delta\omega=(\Delta\omega,0)+DW_\omega\Delta\omega$. With $\pfh=\tilde h+f$,
$X_{\tilde h}(\theta,I)=(\nabla h(I),0)$, and
$\Delta\omega=\nabla h(\iota I_\omega+\iota\Delta I)-\nabla h(\iota I_\omega)$
(Lemma~\ref{bp:p4:freq-map}(ii)),
\begin{equation*}
E'=X_\pfh(K_\omega+(0,\iota\Delta I))-DK_\omega(\omega+\Delta\omega)=(P_h,0)+P_f-DW_\omega\,\Delta\omega\,,
\end{equation*}
where, with $a=\iota I_\omega$ and $v=v_\omega(\theta)$,
\begin{equation*}
P_h=\nabla h(a+v+\iota\Delta I)-\nabla h(a+v)-\nabla h(a+\iota\Delta I)+\nabla h(a)\,,\qquad
P_f=X_f(K_\omega+(0,\iota\Delta I))-X_f(K_\omega)\,.
\end{equation*}

\proofstep{The integrable part: a mixed difference}\label{step:p4:le-h}
For $s',t\in[0,1]$, $\abs{a+s'v+t\iota\Delta I-a}\le\frac34\delta_*<\frac{r_0}2$ with $I_\omega\in B$,
so these points lie in $B_{r_0/2}$. There $\nabla h=J'\circ Dh$ with
$J'(\ell)=(\ell(e_i))_i$ of norm at most $1$, so $\abs{D^2\nabla h}\le\abs{D^3h}\le H_3$
(Lemma~\ref{bp:p4:derivative-bounds}(b)). Lemma~\ref{bp:p4:mixed-difference} gives
$\abs{P_h}\le H_3\eta_0\pfL\abs{\Delta\omega}$.

\proofstep{The perturbation}\label{step:p4:le-f}
Let $\theta\in T_{3s/32}\subset T_{s/4}$. The segment $t\mapsto K_\omega(\theta)+t(0,\iota\Delta I)$
stays within $\frac34\delta_*$ of $K^0_\omega(\theta)$, hence in $\pfD$
(Lemma~\ref{bp:p4:unperturbed}(g)), where $\abs{DX_f}\le F_2\varepsilon$. So
$\abs{P_f}\le F_2\varepsilon\pfL\abs{\Delta\omega}$.

\proofstep{The transport term}\label{step:p4:le-transport}
By Lemma~\ref{bp:p1:cauchy-strip} with loss $\frac s{32}$ from $T_{s/8}$,
$\abs{DW_\omega}\le32n\eta_0/s$ on $T_{3s/32}$.

\proofstep{Sum}\label{step:p4:le-sum}
Since $\gamma\le1$, $\varepsilon\le\varepsilon\gamma^{-2}$, and the three bounds add up to
$\bigl(\pfL H_3A_0+\pfL F_2+32nA_0/s\bigr)\varepsilon\gamma^{-2}\abs{\Delta\omega}=A_1\varepsilon\gamma^{-2}\abs{\Delta\omega}$.
\end{proof}

\begin{theorem}[LIP]\label{bp:p4:lip}
Assume~\eqref{bp:p4:SA}. For all $\omega,\omega'\in\Omega_\gamma$ and all $\theta\in T_{s/32}$, in
particular all real $\theta$,
\begin{equation*}
\abs{W_{\omega'}(\theta)-W_\omega(\theta)}\le A_3\,\varepsilon\gamma^{-4}\abs{\omega'-\omega}\,.
\end{equation*}
\end{theorem}

\begin{proof}
\proofstep{Far pairs}\label{step:p4:lip-far}
If $\abs{\Delta\omega}>\lambda_*\gamma^2$, then by Lemma~\ref{bp:p4:normalize}
$\abs{W_{\omega'}-W_\omega}\le2\eta_0=2A_0\varepsilon\gamma^{-2}<(2A_0/\lambda_*)\varepsilon\gamma^{-4}\abs{\Delta\omega}$.
From now on $\abs{\Delta\omega}\le\lambda_*\gamma^2$; let $K'$ and $E'$ be as in
Lemma~\ref{bp:p4:lip-error} and $e'=A_1\varepsilon\gamma^{-2}\abs{\Delta\omega}$, so that
$e'\le A_1\lambda_*\varepsilon$.

\proofstep{Re-application of G1 at $\omega'$}\label{step:p4:lip-g1}
Apply Theorem~\ref{bp:p0:g1} with the second tuple of Definition~\ref{bp:p4:parameters}
($\rho=\frac{3s}{32}$, $\rho_\infty=\frac s{16}$), $\mathfrak B=\pfD$, $h=\pfh$, $K=K'$, the
frequency $\omega'$, the same $\gamma$, and $e=e'$.
\begin{itemize}
\item H1: Lemma~\ref{bp:p4:h1} and (s1).
\item H2--H4 on $T_{3s/32}$: Theorem~\ref{bp:p4:stab} with $I_*=I_{\omega'}\in B$, frequency
  $\omega'$ ($\abs{\omega'}\le W$), $\rho=\frac s8$, $\delta=\frac s{32}$ and $\eta=\eta_0\le\eta_*(\frac s{32})$
  (s3), using $\abs{K'-K^0_{\omega'}}\le\eta_0$ and that $K'$ is primary and real on reals
  (Lemma~\ref{bp:p4:lip-error}).
\item H5: $0<\gamma\le1$ and $\omega'\in\Dio_{\gamma,\tau}$.
\item Error: $\abs{E'}\le e'$ on $T_{3s/32}$ (Lemma~\ref{bp:p4:lip-error}), and by (s4)
  \begin{equation*}
  \frac{C_1^{(2)}e'}{\gamma^4(3s/32)^{4\tau}}\le\frac{C_1^{(2)}A_1\lambda_*\varepsilon}{\gamma^4(3s/32)^{4\tau}}\le\frac12\,.
  \end{equation*}
\end{itemize}
G1 yields $\tilde K$, holomorphic on $T_{s/16}$, primary, real on reals, invariant with frequency
$\omega'$ on $T_{s/16}$, with $B_{\mu/2}(\tilde K(\theta))\subset\pfD$,
$\abs{D\tilde K}\le\sigma_K-\mu/2$, and $\abs{\tilde K-K'}\le d_1:=C_2^{(2)}e'/(\gamma^2(3s/32)^{2\tau})$
on $T_{s/16}$.

\proofstep{Local uniqueness at $\omega'$}\label{step:p4:lip-g2}
Apply Theorem~\ref{bp:p0:g2} with the tuple $\rho=\frac s{16}$ of Definition~\ref{bp:p4:parameters},
$\mathfrak B=\pfD$, $h=\pfh$, $K_1=K_{\omega'}$, $K_2=\tilde K$, the frequency $\omega'$ and $\gamma$.
\begin{itemize}
\item H1: Lemma~\ref{bp:p4:h1}.
\item $K_1$: H2--H4 on $T_{s/16}$ by Theorem~\ref{bp:p4:stab} ($I_*=I_{\omega'}$, $\rho=\frac s8$,
  $\delta=\frac s{16}$, $\eta=\eta_0\le\eta_*(\frac s{16})$); primary and real on reals, $E_1=0$ on
  $T_{s/16}$, $\gamma>0$ and $\omega'\in\Dio_{\gamma,\tau}$ (Lemma~\ref{bp:p4:normalize}).
\item $K_2$: holomorphic on $T_{s/16}$, primary, real on reals, $E_2=0$, and
  $K_2(T_{s/16})\subset\pfD$ as centres of the balls $B_{\mu/2}(\tilde K(\theta))\subset\pfD$.
\item Closeness on $T_{s/16}$:
  $\abs{K_{\omega'}-\tilde K}\le\abs{W_{\omega'}}+\abs{K^0_{\omega'}-K'}+\abs{K'-\tilde K}\le2\eta_0+d_1\le\bigl(2A_0+C_2^{(2)}A_1\lambda_*(3s/32)^{-2\tau}\bigr)\varepsilon\gamma^{-2}\le c_G\gamma^2(s/16)^{2\tau}$
  by (s5).
\end{itemize}
G2 gives $\varphi\in\R^n$ with $\tilde K(\theta)=K_{\omega'}(\theta+\iota\varphi)$ for $\theta\in T_{s/32}$.

\proofstep{The phase is controlled by normalization}\label{step:p4:lip-phase}
For real $x$, $u_{\tilde K}(\iota x)=\iota\varphi+u_{\omega'}(\iota(x+\varphi))$; averaging with
Lemma~\ref{bp:p4:periodic-average} and $\avg{u_{\omega'}}=0$ gives
$\operatorname{ph}(\tilde K)=\varphi$. The angle components of $K'$ and $K_\omega$ coincide, so
$u_{\tilde K}-u_\omega$ is the angle component of $\tilde K-K'$; since $\avg{u_\omega}=0$,
$\varphi=\int_{[0,1]^n}(u_{\tilde K}-u_\omega)(\iota x)\,dx$ and $\abs\varphi\le d_1$.

\proofstep{Conclusion for close pairs}\label{step:p4:lip-close}
Let $\theta\in T_{s/32}$. Then $K_{\omega'}(\theta)=\tilde K(\theta-\iota\varphi)$, and
$W_{\omega'}-W_\omega=K_{\omega'}-K'$ because $K^0_{\omega'}-K^0_\omega=(0,\iota\Delta I)$. Hence
\begin{equation*}
\abs{W_{\omega'}(\theta)-W_\omega(\theta)}\le\abs{\tilde K(\theta-\iota\varphi)-\tilde K(\theta)}+\abs{\tilde K(\theta)-K'(\theta)}\le\sigma_K\abs\varphi+d_1\le(\sigma_K+1)d_1\,,
\end{equation*}
by Lemma~\ref{bp:p1:increments} on the real segment from $\theta$ to $\theta-\iota\varphi$, which lies
in $T_{s/32}$, with $\abs{D\tilde K}\le\sigma_K$ on $T_{s/16}$. Finally
$(\sigma_K+1)d_1=(\sigma_K+1)C_2^{(2)}A_1(3s/32)^{-2\tau}\varepsilon\gamma^{-4}\abs{\Delta\omega}=A_2\varepsilon\gamma^{-4}\abs{\Delta\omega}$,
and $A_2,2A_0/\lambda_*\le A_3$.
\end{proof}

%======================================================================
\subsection{The real tori}\label{bp:p4:sec-real}

The tori $K_\omega$ live on complex strips; the cover of HEAD consists of real maps
$\R^n\to\R^n\times\R^n$ with the six properties of Definition~\ref{bp:p0:torus}. We restrict to
real points, check analyticity, primarity and the embedding property from the smallness of
$DW_\omega$, derive the Lagrangian property from invariance and nonresonance, and obtain the
Hamiltonian equations from the invariance equation.

\begin{lemma}[Restriction to real points]\label{bp:p4:real-restriction}
Let $K$ be holomorphic on $T_\rho$ and real on reals, and define $k:\R^n\to\R^n\times\R^n$ by
$\iota k(x)=K(\iota x)$ (componentwise real parts). Then:
\begin{enumerate}[label=(\alph*)]
\item $k$ is real-analytic on $\R^n$ (\Lean{AnalyticOnNhd} $\R$ $k$ \Lean{univ});
\item $k$ is differentiable and $\iota(Dk(x)v)=DK(\iota x)(\iota v)$ for $x,v\in\R^n$; hence
  $\abs{Dk(x)}\le\abs{DK(\iota x)}$;
\item if $K$ is primary on $T_\rho$, then $k(x+m)=k(x)+(m,0)$ for $m\in\Z^n$.
\end{enumerate}
\end{lemma}

\begin{proof}
\proofstep{Analyticity}\label{step:p4:rr-analytic}
A function holomorphic on an open subset of $\C^n$ is complex-analytic there
(Mathlib: \Lean{DifferentiableOn.analyticOnNhd}); restricting scalars and composing with the
$\R$-linear map $\iota$ makes $x\mapsto K(\iota x)$ real-analytic on $\R^n$, and $k$ is its
composition with the continuous $\R$-linear map taking real parts.

\proofstep{Derivative}\label{step:p4:rr-derivative}
By the chain rule $D(K\circ\iota)(x)v=DK(\iota x)(\iota v)$. Since $K\circ\iota$ takes values in
$\iota(\R^{2n})$, its difference quotients are real and so is their limit; hence
$D(K\circ\iota)(x)v=\iota(Dk(x)v)$. The norm bound follows because $\iota$ preserves the max norm.
Item (c) is the primarity identity at the real point $\iota x$.
\end{proof}

\begin{lemma}[Real and complex Hamiltonian vector fields]\label{bp:p4:hamiltonian-real}
Assume $\varepsilon\le1$ and let $z\in\R^n\times\R^n$ with $\iota z\in\pfD$. Then
$H=\Lean{realHamiltonian}\,h\,f$ is differentiable at $z$ and
$\iota\bigl(X_H(z)\bigr)=X_\pfh(\iota z)$, where $X_H$ is the frozen \Lean{hamiltonianVectorField}.
\end{lemma}

\begin{proof}
\proofstep{The real Hamiltonian near $z$}\label{step:p4:hr-proof}
By definition $H(y)=\operatorname{Re}\pfh(\iota y)$ for every real $y$. The set
$\{y:\iota y\in\pfD\}$ is an open neighbourhood of $z$ on which $\varphi(y)=\pfh(\iota y)$ is
real-valued (Lemma~\ref{bp:p4:h1}) and $\R$-differentiable with
$D\varphi(y)e=D\pfh(\iota y)(\iota e)$, a real number by the argument of
Step~\ref{step:p4:fm-gradient}. So $H=\varphi$ near $z$ and $DH(z)e=D\pfh(\iota z)(\iota e)$.
Comparing the frozen formulas componentwise,
$X_H(z)=\bigl((DH(z)(0,e_i))_i,(-DH(z)(e_i,0))_i\bigr)$ is the real point whose image under $\iota$ is
$X_\pfh(\iota z)$.
\end{proof}

\begin{lemma}[Invariant tori with nonresonant frequency are Lagrangian]\label{bp:p4:lagrangian}
Assume $\varepsilon\le1$. Let $\rho>0$, $\gamma>0$, $\omega\in\Dio_{\gamma,\tau}$, and let $K$ be
holomorphic on $T_\rho$, primary on $T_\rho$, with $K(T_\rho)\subset\pfD$, $DK$ bounded on $T_\rho$,
and $E=0$ at $\omega$ on $T_\rho$. Then $\Lambda(\theta):=DK(\theta)^\top\Omega_0DK(\theta)=0$ for
$\theta\in T_\rho$.
\end{lemma}

\source{GJLV05 Lemma~19(b),(d), pp.~889--891 (``In particular, $L=0$ if $e=0$''), for vector
fields with a primitive of the symplectic form; we give a direct proof in our coordinates.}
\sourcegap{HL19 \S2.2, p.~7, eq.~(2.3), asserts that quasi-periodicity implies isotropy (``It is
well known \dots\ see [6, 41]'') without proof. The proof below supplies it.}

\begin{proof}
\proofstep{Regularity and periodicity}\label{step:p4:lag-regular}
Write $K=(\theta+u,v)$. By Step~\ref{step:p4:tr-periodic}, $u$, $v$ and $DK$ are $\Z^n$-periodic
on $T_\rho$; $\Lambda$ is holomorphic, periodic and bounded on $T_\rho$. On $T_{\rho/2}$ the
functions $u,v,Du$ are bounded by Lemma~\ref{bp:p1:periodic-bounded}.

\proofstep{$\Lambda$ is a sum of derivatives}\label{step:p4:lag-exact}
By Lemma~\ref{bp:p4:matrix-norms}(iv),
$\Lambda_{ij}=\sum_m(\partial_iv_m\,\partial_jK_{1,m}-\partial_iK_{1,m}\,\partial_jv_m)$ with
$\partial_jK_{1,m}=\delta_{jm}+\partial_ju_m$. Using $\partial_i\partial_ju=\partial_j\partial_iu$,
\begin{equation*}
\Lambda_{ij}=\partial_iv_j-\partial_jv_i+\partial_i(v\cdot\partial_ju)-\partial_j(v\cdot\partial_iu)\,.
\end{equation*}
Each term is a partial derivative of a bounded periodic holomorphic function on $T_{\rho/2}$, so
$\avg\Lambda=0$ by Lemma~\ref{bp:p1:averages}.

\proofstep{$\Lambda$ is constant along $\omega$}\label{step:p4:lag-transport}
Differentiating $X_\pfh(K(\theta))=DK(\theta)\omega$ in the direction $e_i$ and using the
symmetry of $D^2K$ gives $DX_\pfh(K)\partial_iK=\partial_\omega\partial_iK$. By
Step~\ref{step:p4:db-vectorfield}, $DX_\pfh(z)=-\Omega_0S(z)$, where $S(z)$ is the complex Hessian
matrix of $\pfh$ in the coordinates $(\theta,I)$, which is symmetric. Hence, by
Lemma~\ref{bp:p4:matrix-norms}(iv),
\begin{equation*}
\partial_\omega\Lambda_{ij}=-(\Omega_0S\partial_iK)^\top\Omega_0\partial_jK-\partial_iK^\top\Omega_0\Omega_0S\partial_jK
=-\partial_iK^\top S^\top\partial_jK+\partial_iK^\top S\partial_jK=0\,.
\end{equation*}

\proofstep{Uniqueness for the cohomological equation}\label{step:p4:lag-unique}
Apply Lemma~\ref{bp:p1:cohomological} on $T_{\rho/2}$ with $g=0$ and loss $\rho/4$, entrywise:
there is exactly one holomorphic periodic $w$ on $T_{\rho/4}$ with $\avg w=0$ and $Dw\,\omega=0$.
Both $w=0$ and $w=\Lambda_{ij}$ qualify (Steps~\ref{step:p4:lag-exact}
and~\ref{step:p4:lag-transport}), so $\Lambda=0$ on $T_{\rho/4}$, and on $T_\rho$ by the identity
theorem.
\end{proof}

\begin{theorem}[The real tori are KAM tori]\label{bp:p4:real-torus}
Assume~\eqref{bp:p4:SA} and let $\omega\in\Omega_\gamma$. Let $k_\omega$ be the restriction of
$K_\omega$ to real points (Lemma~\ref{bp:p4:real-restriction}). Then:
\begin{enumerate}[label=(\roman*)]
\item $k_\omega$ is an invariant Lagrangian torus of $H=\Lean{realHamiltonian}\,h\,f$ with
  frequency $\omega$ (Definition~\ref{bp:p0:torus});
\item for every $x\in\R^n$, $\abs{k_\omega(x)_2-I_\omega}<r_0$, with $I_\omega\in B$;
\item $I_\omega\in B$ and $\pfg(I_\omega)=\omega$;
\item $\abs{k_\omega(x)-(x,I_\omega)}\le\eta_0$ for every $x\in\R^n$.
\end{enumerate}
Consequently $k_\omega(\R^n)\subset\mathcal K(\eta_0)=\Lean{kamCover}\,h\,f\,B\,r_0\,\tau\,\eta_0$,
with witnesses $K=k_\omega$, $\omega$, $\gamma$ and $I_0=I_\omega$.
\end{theorem}

\begin{proof}
\proofstep{Analyticity and primarity}\label{step:p4:rt-analytic}
Lemma~\ref{bp:p4:real-restriction}(a),(c) and Lemma~\ref{bp:p4:normalize}.

\proofstep{Small angle displacement}\label{step:p4:rt-small}
Let $w(x)=k_\omega(x)-(x,I_\omega)$, the real restriction of $W_\omega$, and $\tilde u=w_1$. By
Lemma~\ref{bp:p1:cauchy-strip} from $T_{s/8}$ with loss $\frac s{16}$ and (s3),
$\abs{DW_\omega}\le16n\eta_0/s\le\frac12$ on $T_{s/16}\supset\iota\R^n$; by
Lemma~\ref{bp:p4:real-restriction}(b), $\abs{D\tilde u(x)}\le\frac12$ for every $x\in\R^n$, and
$\tilde u$ is $\Z^n$-periodic.

\proofstep{Immersion}\label{step:p4:rt-immersion}
$Dk_\omega(x)v$ has angle component $v+D\tilde u(x)v$, of max norm at least $\frac12\abs v$, so
$Dk_\omega(x)$ is injective.

\proofstep{Embedding modulo $\Z^n$}\label{step:p4:rt-embedding}
If $k_\omega(x')=k_\omega(x)+(m,0)$ with $m\in\Z^n$, the angle components give
$x'+\tilde u(x')=x+m+\tilde u(x)=x+m+\tilde u(x+m)$. Since $\R^n$ is convex and
$\abs{D\tilde u}\le\frac12$, $\abs{x'-(x+m)}=\abs{\tilde u(x+m)-\tilde u(x')}\le\frac12\abs{x'-(x+m)}$,
so $x'=x+m$.

\proofstep{Lagrangian}\label{step:p4:rt-lagrangian}
Lemma~\ref{bp:p4:lagrangian} applies to $K_\omega$ on $T_{s/8}$: it is holomorphic and primary,
$K_\omega(T_{s/8})\subset\pfD$ and $\abs{DK_\omega}\le\sigma_K$ (Lemma~\ref{bp:p4:normalize}),
$E=0$, $\gamma>0$, $\omega\in\Dio_{\gamma,\tau}$, and $\varepsilon\le1$ by (s1). For real $x,u,v$,
Lemma~\ref{bp:p4:matrix-norms}(iv) and Lemma~\ref{bp:p4:real-restriction}(b) give
$\Lean{symplecticForm}(Dk_\omega(x)u,Dk_\omega(x)v)=(\iota u)^\top\Lambda(\iota x)(\iota v)=0$.

\proofstep{Hamilton's equations}\label{step:p4:rt-ode}
Fix $x\in\R^n$, $t\in\R$, and put $y=x+t\omega$. Then $\iota k_\omega(y)=K_\omega(\iota y)\in\pfD$,
so by Lemma~\ref{bp:p4:hamiltonian-real}, invariance at $\iota y\in T_{s/8}$, and
Lemma~\ref{bp:p4:real-restriction}(b),
$\iota X_H(k_\omega(y))=X_\pfh(K_\omega(\iota y))=DK_\omega(\iota y)(\iota\omega)=\iota\bigl(Dk_\omega(y)\omega\bigr)$.
By the chain rule, $s\mapsto k_\omega(x+s\omega)$ has derivative $Dk_\omega(y)\omega=X_H(k_\omega(y))$
at $s=t$, which is clause (vi) of Definition~\ref{bp:p0:torus}.

\proofstep{Cover clauses}\label{step:p4:rt-cover}
$\abs{k_\omega(x)-(x,I_\omega)}=\abs{W_\omega(\iota x)}\le\eta_0$ (Lemma~\ref{bp:p4:normalize}),
which is (iv); in particular $\abs{k_\omega(x)_2-I_\omega}\le\eta_0\le\mu<r_0$ by (s3), which is (ii)
with $J=I_\omega\in B$; (iii) is Lemma~\ref{bp:p4:inverse-freq}. Together with $\gamma>0$ and
$\omega\in\Dio_{\gamma,\tau}$ these are all the clauses of Definition~\ref{bp:p0:cover} with
$\delta=\eta_0$.
\end{proof}

%======================================================================
\subsection{The measure of the complement (HEAD)}\label{bp:p4:sec-measure}

We parameterize the tori by their actions: for a good action $I$ the torus
$\theta\mapsto k_{\pfg(I)}(\theta)$ passes within $\eta_0$ of $(\theta,I)$, and by
Theorem~\ref{bp:p4:lip} the displacement $(\theta,I)\mapsto k_{\pfg(I)}(\theta)-(\theta,I)$ is
$\frac12$-Lipschitz once $\varepsilon\lesssim\gamma^4$. We extend this displacement by zero near
the boundary of $B$, then to all of $\R^{2n}$ by McShane's formula, so that the identity plus the
extension is a lipeomorphism of $\R^{2n}$ fixing the boundary layer. A point of
$[0,1)^n\times B$ off the tori is then the image of a point whose action is bad: too close to the
boundary, or with a frequency that is not $(\gamma,\tau)$-Diophantine. Both sets of bad actions have
measure $O(\gamma)$.
\source{Pos01 Prop.~1--2, p.~11, and App.~B, p.~31; BC17 \S6.6, Lemma~6 and footnote~32, pp.~17--18
(ledger D-MEAS-1). We work in action coordinates rather than frequency coordinates; see
Remark~\ref{bp:p4:omega-gamma} and the ledger candidates.}

\begin{lemma}[Lipschitz images]\label{bp:p4:lip-image}
Let $N\ge1$, $S\subset\R^N$ any set, and $F:S\to\R^N$ with $\abs{F(x)-F(y)}\le\lambda\abs{x-y}$
for the max norm. Then $\meas F(S)\le\lambda^N\meas S$ (outer measures). The same holds on
$\R^n\times\R^n$ with the max norm of the two components and the product Lebesgue measure.
\end{lemma}

\begin{proof}
\proofstep{Hausdorff measure}\label{step:p4:li-proof}
For the max norm on $\R^N$ the $N$-dimensional Hausdorff measure equals Lebesgue measure
(Mathlib: \nolinkurl{MeasureTheory.hausdorffMeasure_pi_real}), and a $\lambda$-Lipschitz map on any
set multiplies $\mu_H^N$ by at most $\lambda^N$
(\nolinkurl{LipschitzOnWith.hausdorffMeasure_image_le}).
The identification $\R^n\times\R^n\cong\R^{2n}$ is an isometry for the max norms and carries the
product of Lebesgue measures to Lebesgue measure.
\end{proof}
\source{BC17, footnote~32, p.~18: ``If $A\subseteq\R^n$ is a measurable set and $f:A\to\R^n$ is a
Lipschitz map, then $\meas f(A)\le\abs f^n_{\mathrm{Lip},A}\meas A$.''}

\begin{lemma}[Actions with non-Diophantine frequency]\label{bp:p4:preimage-measure}
Let $C_3$ be the constant of Theorem~\ref{bp:p0:g3} for the bounded set $\Omega=\pfg(B)$
(Lemma~\ref{bp:p4:freq-map}(v)) and $C_3^+=\max(C_3,0)$. For every $\alpha>0$,
\begin{equation*}
\meas\{I\in B:\pfg(I)\notin\Dio_{\alpha,\tau}\}\le\pfL^n\,C_3^+\,\alpha\,.
\end{equation*}
\end{lemma}

\begin{proof}
\proofstep{Image under the inverse frequency map}\label{step:p4:pm-proof}
The set is the image of $\{\omega\in\pfg(B):\omega\notin\Dio_{\alpha,\tau}\}$ under the inverse of
$\pfg$, which is $\pfL$-Lipschitz on $\pfg(\bar B)$ (Lemma~\ref{bp:p4:inverse-freq}). Apply
Lemma~\ref{bp:p4:lip-image} with $N=n$ and Theorem~\ref{bp:p0:g3}.
\end{proof}

\begin{lemma}[McShane extension with a sup bound]\label{bp:p4:mcshane}
Let $X$ be a metric space, $S\subset X$, $\lambda,\beta\ge0$, and $\phi:S\to\R^m$ with
$\abs{\phi(x)-\phi(y)}\le\lambda d(x,y)$ and $\abs{\phi(x)}\le\beta$ (max norm) for $x,y\in S$.
Then there is $\Phi:X\to\R^m$ with $\Phi=\phi$ on $S$, $\abs{\Phi(x)-\Phi(y)}\le\lambda d(x,y)$
and $\abs\Phi\le\beta$ on $X$.
\end{lemma}

\begin{proof}
\proofstep{One component}\label{step:p4:mc-component}
If $S=\emptyset$ take $\Phi=0$. Otherwise, for each component $\phi_j$ put
$U_j(x)=\inf_{z\in S}\bigl(\phi_j(z)+\lambda d(x,z)\bigr)$. It is finite, since every term is at least
$\phi_j(z_0)-\lambda d(x,z_0)$ for a fixed $z_0\in S$; it is $\lambda$-Lipschitz as an infimum of
$\lambda$-Lipschitz functions; and $U_j=\phi_j$ on $S$ because $\phi_j(z)+\lambda d(x,z)\ge\phi_j(x)$
for $x,z\in S$, with equality at $z=x$. (This is the formula of Pos01 App.~B with $\sup$ and $\inf$
exchanged; Mathlib: \Lean{LipschitzOnWith.extend\_real}.)

\proofstep{Clamping}\label{step:p4:mc-clamp}
Let $\Phi_j=\max(-\beta,\min(\beta,U_j))$. Clamping is $1$-Lipschitz and fixes $[-\beta,\beta]$, so
$\Phi_j$ is $\lambda$-Lipschitz, extends $\phi_j$, and $\abs{\Phi_j}\le\beta$. In the max norm,
$\Phi=(\Phi_j)_j$ is $\lambda$-Lipschitz because each component is.
\end{proof}

\begin{lemma}[Lipeomorphisms]\label{bp:p4:lipeo}
Let $G:\R^N\to\R^N$ be $\eta$-Lipschitz with $\eta<1$ and $\Psi=\mathrm{id}+G$. Then $\Psi$ is a
bijection, $(1-\eta)\abs{x-y}\le\abs{\Psi(x)-\Psi(y)}\le(1+\eta)\abs{x-y}$, and if $G=0$ on a set
$Z'$ then $\Psi$ maps $\R^N\setminus Z'$ onto itself.
\end{lemma}

\begin{proof}
\proofstep{Bi-Lipschitz bounds and injectivity}\label{step:p4:lipeo-inj}
The two inequalities are the triangle inequality; the lower one gives injectivity.

\proofstep{Surjectivity}\label{step:p4:lipeo-onto}
For $y\in\R^N$, the map $x\mapsto y-G(x)$ is a contraction of the complete space $\R^N$, so it has
a fixed point $x$, and $\Psi(x)=y$ (Banach fixed point theorem; Mathlib:
\Lean{ContractingWith.fixedPoint}).

\proofstep{The fixed region}\label{step:p4:lipeo-fixed}
If $x\notin Z'$ and $\Psi(x)\in Z'$, then $\Psi(\Psi(x))=\Psi(x)$, so $\Psi(x)=x\notin Z'$ by
injectivity, a contradiction. If $y\notin Z'$ and $y=\Psi(x)$ with $x\in Z'$, then $y=x\in Z'$, a
contradiction. So $\Psi$ maps $\R^N\setminus Z'$ into and onto itself.
\end{proof}
\source{Pos01 Prop.~1, p.~11 (``Hence $\bar\varphi$ is a lipeomorphism on $\R^n$''); BC17 footnote~28,
p.~17.}

\begin{definition}[Good actions and the displacement]\label{bp:p4:good-actions}
Assume~\eqref{bp:p4:SA}. Put
\begin{gather*}
Z=\{I\in\R^n:\infd(I,B^c)\le\gamma\}\,,\qquad
A_\gamma=\{I\in B:\infd(I,B^c)\ge2\gamma,\ \pfg(I)\in\Dio_{\gamma,\tau}\}\,,\\
G_0(\theta,I)=k_{\pfg(I)}(\theta)-(\theta,I)\quad\text{on }\R^n\times A_\gamma\,,\qquad G_0=0\quad\text{on }\R^n\times Z\,.
\end{gather*}
For $I\in A_\gamma$, $\pfg(I)\in\Omega_\gamma$ and $I_{\pfg(I)}=I$ (Lemma~\ref{bp:p4:inverse-freq}),
so $G_0(\theta,I)$ is the real restriction of $W_{\pfg(I)}$ at $\theta$.
\end{definition}

\begin{lemma}[The displacement is small and $\tfrac12$-Lipschitz]\label{bp:p4:map-lipschitz}
Assume~\eqref{bp:p4:SA}. The sets $A_\gamma$ and $Z$ are disjoint, $\abs{G_0}\le\eta_0$, and
$G_0$ is $\frac12$-Lipschitz on $\R^n\times(A_\gamma\cup Z)$ for the max norm.
\end{lemma}

\begin{proof}
\proofstep{Disjointness and size}\label{step:p4:ml-size}
Points of $A_\gamma$ have $\infd(I,B^c)\ge2\gamma>\gamma$. On $\R^n\times A_\gamma$,
$\abs{G_0}\le\sup\abs{W_{\pfg(I)}}\le\eta_0$ by Lemma~\ref{bp:p4:normalize}.

\proofstep{Two good actions}\label{step:p4:ml-good}
Let $I,J\in A_\gamma$ and $\theta,\theta'\in\R^n$. By Step~\ref{step:p4:rt-small} the real
restriction of $W_\omega$ has derivative at most $16n\eta_0/s=A_4\varepsilon\gamma^{-2}$ on the
convex set $\R^n$; by Theorem~\ref{bp:p4:lip} and Lemma~\ref{bp:p4:freq-map}(v),
\begin{align*}
\abs{G_0(\theta,I)-G_0(\theta',J)}&\le\abs{W_{\pfg(I)}(\theta)-W_{\pfg(I)}(\theta')}+\abs{W_{\pfg(I)}(\theta')-W_{\pfg(J)}(\theta')}\\
&\le A_4\varepsilon\gamma^{-2}\abs{\theta-\theta'}+A_3\varepsilon\gamma^{-4}\Lambda\abs{I-J}\le\tfrac14\abs{(\theta,I)-(\theta',J)}\,,
\end{align*}
by (s6) and $\gamma\le1$.

\proofstep{A good action and a boundary action}\label{step:p4:ml-mixed}
Let $I\in A_\gamma$ and $J\in Z$. Since $\infd(\cdot,B^c)$ is $1$-Lipschitz,
$\abs{I-J}\ge2\gamma-\gamma=\gamma$, and by (s7)
$\abs{G_0(\theta,I)-G_0(\theta',J)}\le\eta_0\le\frac12\gamma\le\frac12\abs{(\theta,I)-(\theta',J)}$. Two
boundary actions give $0$.
\end{proof}

\begin{theorem}[Measure of the complement at a given $\gamma$]\label{bp:p4:main-measure}
Assume~\eqref{bp:p4:SA} and $B\ne\emptyset$, and put
$C_{\mathrm{meas}}=2C_B^++(9/2)^n\bigl(2C_B^++\pfL^nC_3^+\bigr)$. Then
\begin{equation*}
\meas\bigl(F_B\setminus\mathcal K(\eta_0)\bigr)\le C_{\mathrm{meas}}\,\gamma\,,\qquad
\mathcal K(\eta_0)=\Lean{kamCover}\,h\,f\,B\,r_0\,\tau\,\eta_0\,.
\end{equation*}
\end{theorem}

\begin{proof}
\proofstep{Extension and lipeomorphism}\label{step:p4:mm-extension}
By Lemmas~\ref{bp:p4:map-lipschitz} and~\ref{bp:p4:mcshane} (with $X=\R^n\times\R^n$, max norm,
$S=\R^n\times(A_\gamma\cup Z)$, $\lambda=\frac12$, $\beta=\eta_0$) there is
$\bar G:\R^{2n}\to\R^{2n}$, $\frac12$-Lipschitz, with $\bar G=G_0$ on $S$ and $\abs{\bar G}\le\eta_0\le\frac12$.
By Lemma~\ref{bp:p4:lipeo}, $\Psi=\mathrm{id}+\bar G$ is a bijection of $\R^{2n}$, it is
$\frac32$-Lipschitz, and it maps $\R^n\times Z^c$ onto itself because $\bar G=0$ on $\R^n\times Z$.

\proofstep{Good actions lie on tori}\label{step:p4:mm-tori}
For $I\in A_\gamma$ and $\theta\in\R^n$, $\Psi(\theta,I)=k_{\pfg(I)}(\theta)\in\mathcal K(\eta_0)$ by
Theorem~\ref{bp:p4:real-torus} with $\omega=\pfg(I)\in\Omega_\gamma$.

\proofstep{The complement is the image of bad actions}\label{step:p4:mm-complement}
Let $y=(y_1,y_2)\in F_B\setminus\mathcal K(\eta_0)$. If $y_2\in Z$, then $y\in[0,1)^n\times(B\cap Z)$.
Otherwise $y=\Psi(x)$ for some $x\in\R^n\times Z^c$, and $x\notin\R^n\times A_\gamma$ by
Step~\ref{step:p4:mm-tori}. Since $\abs{x-y}=\abs{\bar G(x)}\le\frac12$, $x_1$ lies in
$Q=[-\frac12,\frac32]^n$. Hence
\begin{equation*}
F_B\setminus\mathcal K(\eta_0)\subset\bigl([0,1)^n\times(B\cap Z)\bigr)\cup\Psi\bigl(Q\times(Z^c\setminus A_\gamma)\bigr)\,.
\end{equation*}

\proofstep{Bad actions}\label{step:p4:mm-bad}
Let $Y=\{I\in B:\infd(I,B^c)<2\gamma\}$. If $I\notin B$ then $\infd(I,B^c)=0$, so $Z^c\subset B$.
Therefore $B\cap Z\subset Y$ and $Z^c\setminus A_\gamma\subset Y\cup\{I\in B:\pfg(I)\notin\Dio_{\gamma,\tau}\}$.
By the boundary-layer hypothesis with $\delta=2\gamma$, $\meas Y\le2C_B^+\gamma$, and by
Lemma~\ref{bp:p4:preimage-measure} the second set has measure at most $\pfL^nC_3^+\gamma$.

\proofstep{Measure count}\label{step:p4:mm-count}
The product measure of $[0,1)^n\times(B\cap Z)$ is $\meas(B\cap Z)\le2C_B^+\gamma$. By
Lemma~\ref{bp:p4:lip-image} with $N=2n$ and $\lambda=\frac32$,
\begin{equation*}
\meas\Psi\bigl(Q\times(Z^c\setminus A_\gamma)\bigr)\le(\tfrac32)^{2n}\,2^n\,\meas(Z^c\setminus A_\gamma)\le(\tfrac92)^n\bigl(2C_B^++\pfL^nC_3^+\bigr)\gamma\,.
\end{equation*}
Subadditivity gives the claim.
\end{proof}

\begin{lemma}[Monotonicity of the cover]\label{bp:p4:cover-mono}
If $B'\subset B$ and $\delta'\le\delta$, then
$\Lean{kamCover}\,h\,f\,B'\,r_0\,\tau\,\delta'\subset\Lean{kamCover}\,h\,f\,B\,r_0\,\tau\,\delta$.
\end{lemma}

\begin{proof}
\proofstep{Same witnesses}\label{step:p4:cm-proof}
Keep the witnesses $K,\omega,\gamma,I_0$ of Definition~\ref{bp:p0:cover}. The clauses $\gamma>0$,
Diophantine, invariant Lagrangian torus of $\Lean{realHamiltonian}\,h\,f$ (which does not involve
$B$), $\pfg(I_0)=\omega$ and $z\in K(\R^n)$ are unchanged. A point $J\in B'$ with
$\abs{K(\theta)_2-J}<r_0$ lies in $B$; $I_0\in B'\subset B$; and
$\abs{K(\theta)-(\theta,I_0)}\le\delta'\le\delta$.
\end{proof}

\begin{proof}[Proof of Theorem~\ref{bp:p0:head}]
\proofstep{Empty action domain}\label{step:p4:head-empty}
If $B=\emptyset$, then $F_B=\emptyset$ and $C=\varepsilon_0=1$ work. Let $B\ne\emptyset$.

\proofstep{Constants}\label{step:p4:head-constants}
With $c_*$, $A_0$ (Definition~\ref{bp:p4:parameters}) and $C_{\mathrm{meas}}$
(Theorem~\ref{bp:p4:main-measure}), put
\begin{equation*}
C=\max\bigl(1,\ A_0c_*^{1/2},\ C_{\mathrm{meas}}c_*^{-1/4}\bigr)\,,\qquad\varepsilon_0=c_*\,.
\end{equation*}
Both are positive and depend only on the data $n,\tau,s_0,r_0,B,h$ and the witnesses $L,C_B$,
which are fixed once from the hypotheses. Let $f$ and $\varepsilon\le\varepsilon_0$ satisfy the
hypotheses of HEAD.

\proofstep{Negative $\varepsilon$}\label{step:p4:head-negative}
Since $B\ne\emptyset$, $D(s_0,r_0)\ne\emptyset$, and $\abs{f(z)}\le\varepsilon$ at one of its points
forces $\varepsilon\ge0$.

\proofstep{Positive $\varepsilon$}\label{step:p4:head-positive}
If $0<\varepsilon\le c_*$, put $\gamma=(\varepsilon/c_*)^{1/4}\in(0,1]$. Then $\varepsilon=c_*\gamma^4$,
so~\eqref{bp:p4:SA} holds, and $\eta_0=A_0\varepsilon\gamma^{-2}=A_0c_*^{1/2}\varepsilon^{1/2}\le C\sqrt\varepsilon$.
By Lemma~\ref{bp:p4:cover-mono} and Theorem~\ref{bp:p4:main-measure},
\begin{equation*}
\meas\bigl(F_B\setminus\mathcal K(C\sqrt\varepsilon)\bigr)\le\meas\bigl(F_B\setminus\mathcal K(\eta_0)\bigr)\le C_{\mathrm{meas}}\gamma=C_{\mathrm{meas}}c_*^{-1/4}\varepsilon^{1/4}\le C\varepsilon^{1/4}\,.
\end{equation*}

\proofstep{Vanishing perturbation}\label{step:p4:head-zero}
If $\varepsilon=0$, then~\eqref{bp:p4:SA} holds for every $\gamma\in(0,1]$, with $\eta_0=0=C\sqrt0$.
Theorem~\ref{bp:p4:main-measure} gives $\meas(F_B\setminus\mathcal K(0))\le C_{\mathrm{meas}}\gamma$
for every such $\gamma$, so the measure is $0=C\cdot0^{1/4}$.
\end{proof}

%======================================================================
\subsection{The qualitative form (HEAD-Q)}\label{bp:p4:sec-headq}

HEAD-Q drops the boundary hypothesis. We approximate $B$ from inside by a finite union of open
dyadic cubes, which has the boundary-layer property, apply HEAD there, and absorb the small
remainder $B\setminus B'$ into the allowed fraction $\kappa$.

\begin{lemma}[Inner dyadic approximation]\label{bp:p4:dyadic-approx}
Let $B\subset\R^n$ be open and bounded and $\eta>0$. There is an open set $B'\subset B$ with
$\meas(B\setminus B')\le\eta$ and a constant $C_{B'}$ with
$\meas\{I\in B':\infd(I,B'^c)<\delta\}\le C_{B'}\delta$ for all $\delta>0$.
\end{lemma}

\begin{proof}
\proofstep{Dyadic exhaustion}\label{step:p4:dy-exhaust}
For $j\in\N$ let $\mathcal D_j$ be the set of closed cubes $2^{-j}(m+[0,1]^n)$, $m\in\Z^n$, contained
in $B$; it is finite because $B$ is bounded. Let $E_j=\bigcup\mathcal D_j$, a closed set. Each cube of
$\mathcal D_j$ is the union of $2^n$ cubes of generation $j+1$, so $E_j\subset E_{j+1}$. If $x\in B$
and the ball of radius $r$ around $x$ lies in $B$, a generation-$j$ cube containing $x$ has max-norm
diameter $2^{-j}<r$ for large $j$ and lies in $B$; so $\bigcup_jE_j=B$. By continuity of measure
from below and $\meas B<\infty$, $\meas(B\setminus E_j)\le\eta$ for some $j$.

\proofstep{The open set}\label{step:p4:dy-open}
Let $B'$ be the union of the interiors of the $N$ cubes of $\mathcal D_j$, of side $\ell=2^{-j}$.
It is open, $B'\subset E_j\subset B$, and $E_j\setminus B'$ lies in the union of the cube boundaries, a
null set. Hence $\meas(B\setminus B')\le\eta$.

\proofstep{Boundary layer}\label{step:p4:dy-layer}
Let $I\in B'$ with $\infd(I,B'^c)<\delta$, and let $Q$ be a cube of $\mathcal D_j$ with
$I\in\operatorname{int}Q$. Since $B'^c\subset(\operatorname{int}Q)^c$,
$\infd(I,(\operatorname{int}Q)^c)<\delta$. For $Q=c+[0,\ell]^n$ in the max norm this means that some
coordinate satisfies $I_i-c_i<\delta$ or $c_i+\ell-I_i<\delta$, so $I\in Q\setminus(c+[\delta,\ell-\delta]^n)$,
a set of measure $\ell^n-(\ell-2\delta)_+^n\le2n\delta\ell^{n-1}$. Summing over the cubes,
$C_{B'}=2nN\ell^{n-1}$ works.
\end{proof}

\begin{proof}[Proof of Theorem~\ref{bp:p0:headq}]
\proofstep{Trivial cases}\label{step:p4:hq-trivial}
If $B=\emptyset$, $F_B=\emptyset$ and $\varepsilon_0=1$ works. Let $B\ne\emptyset$; then
$0<\meas B<\infty$ since $B$ is open and bounded, and $\meas F_B=\meas([0,1)^n)\meas B=\meas B$.
Fix $\kappa>0$.

\proofstep{Inner approximation}\label{step:p4:hq-approx}
Lemma~\ref{bp:p4:dyadic-approx} with $\eta=\frac12\kappa\meas B$ gives $B'\subset B$.

\proofstep{HEAD for $B'$}\label{step:p4:hq-head}
The data $(n,\tau,s_0,r_0,B',h)$ satisfy the hypotheses of Theorem~\ref{bp:p0:head}: $B'$ is open,
bounded (as a subset of $B$) and has the boundary-layer property with $C_{B'}$; $B'_{r_0}\subset B_{r_0}$,
so $h$ is holomorphic on $B'_{r_0}$ and real at its real points; and
$\overline{B'}\subset\bar B$, so the nondegeneracy inequality holds on $\overline{B'}$ with the same
$L$. HEAD gives $C'>0$ and $\varepsilon_0'>0$. Put
\begin{equation*}
\varepsilon_0=\min\Bigl(\varepsilon_0',\ \bigl(\kappa/C'\bigr)^2,\ \bigl(\kappa\meas B/(2C')\bigr)^4\Bigr)>0\,.
\end{equation*}

\proofstep{Transfer of the perturbation hypotheses}\label{step:p4:hq-transfer}
Let $f$, $\varepsilon\le\varepsilon_0$ satisfy the hypotheses of HEAD-Q for $B$. Since
$T_{s_0}\times B'_{r_0}\subset T_{s_0}\times B_{r_0}$, they hold for $B'$: holomorphy, periodicity,
reality and $\abs f\le\varepsilon$ restrict, and $\varepsilon\le\varepsilon_0'$. As in
Step~\ref{step:p4:head-negative}, $\varepsilon\ge0$. HEAD gives
$\meas\bigl(F_{B'}\setminus\mathcal K_{B'}(C'\sqrt\varepsilon)\bigr)\le C'\varepsilon^{1/4}\le\frac12\kappa\meas B$,
where $\mathcal K_{B'}$ is the cover built with $B'$.

\proofstep{Comparison of covers}\label{step:p4:hq-cover}
$C'\sqrt\varepsilon\le\kappa$, so $\mathcal K_{B'}(C'\sqrt\varepsilon)\subset\mathcal K_B(\kappa)$ by
Lemma~\ref{bp:p4:cover-mono}. Hence
$F_B\setminus\mathcal K_B(\kappa)\subset\bigl([0,1)^n\times(B\setminus B')\bigr)\cup\bigl(F_{B'}\setminus\mathcal K_{B'}(C'\sqrt\varepsilon)\bigr)$,
and
\begin{equation*}
\meas\bigl(F_B\setminus\mathcal K_B(\kappa)\bigr)\le\meas(B\setminus B')+\tfrac12\kappa\meas B\le\kappa\meas B=\kappa\meas F_B\,.
\end{equation*}
\end{proof}

%======================================================================
\subsection{Applications of G1 and G2, transfers and ledger candidates}\label{bp:p4:sec-notes}

\paragraph{Applications of G1, G2 and G3.} G1 (Theorem~\ref{bp:p0:g1}) is applied twice:
in Step~\ref{step:p4:apply-g1} at $K^0_\omega$ on $T_{s/4}$ (output on $T_{s/8}$), and in
Step~\ref{step:p4:lip-g1} at $K_\omega+(0,\iota\Delta I)$ on $T_{3s/32}$ (output on $T_{s/16}$).
G2 (Theorem~\ref{bp:p0:g2}) is applied twice on $T_{s/16}$: in Step~\ref{step:p4:nu-g2} (uniqueness
of the normalized torus) and in Step~\ref{step:p4:lip-g2} (comparison at $\omega'$). Every
application uses $\mathfrak B=\pfD$, $h=\pfh$ and the parameters of
Definition~\ref{bp:p4:parameters}; each lists the frozen hypotheses with the lemma that checks them.
G3 is applied once, in Lemma~\ref{bp:p4:preimage-measure}.

\paragraph{Transfer requests to Part~\ref{bp:p1}.} The following infrastructure is stated and
proved here and may move to Part~\ref{bp:p1}: Lemma~\ref{bp:p4:matrix-norms}(iv) (the bilinear
form of $\Omega_0$ and \Lean{symplecticForm}; the other items are already in Part~\ref{bp:p1});
Lemma~\ref{bp:p4:torsion-regular} (guarded facts for the Gram inverse, the
frame and the torsion; alternatively Part~\ref{bp:p2}); Lemma~\ref{bp:p4:periodic-average};
Lemma~\ref{bp:p4:real-restriction}; Lemma~\ref{bp:p4:hamiltonian-real};
Lemma~\ref{bp:p4:mixed-difference}; Lemma~\ref{bp:p4:lip-image}; Lemma~\ref{bp:p4:mcshane};
Lemma~\ref{bp:p4:lipeo}; the identity $X_\phi=J\circ D\phi$ with its consequences
(Step~\ref{step:p4:db-vectorfield}); the second-derivative Cauchy estimate on strips obtained by two
applications of Lemma~\ref{bp:p1:cauchy-strip} (Step~\ref{step:p4:stab-cauchy}).

\paragraph{Requests to the orchestrator.}
\begin{enumerate}[label=(R\arabic*)]
\item Lemma~\ref{bp:p1:cohomological} now states uniqueness without a boundedness assumption;
  Lemma~\ref{bp:p4:lagrangian} uses it with $g=0$. Please keep $g=0$ admissible (it is: $\avg0=0$).
\item Lemma~\ref{bp:p1:cauchy-open}: please fix that $D^j$ is the $j$-th Fréchet derivative
  (\Lean{iteratedFDeriv}) with the operator norm for max norms. Part~\ref{bp:p4} applies it to
  scalar functions only and passes to $X_\phi$ through Step~\ref{step:p4:db-vectorfield}.
\item Part~\ref{bp:p0} uses $B$ both for the action domain and for the Gram inverse
  (Definition~\ref{bp:p0:stripdata}); Part~\ref{bp:p4} writes $G_K^{-1}$ for the latter. A common
  renaming would help the graph extraction.
\item \texttt{docs/statement.md} lists $B$ as nonempty; the frozen Lean statements do not, and the
  proofs above treat $B=\emptyset$ separately. No change is needed.
\item The strip losses of the route verdict (\texttt{docs/PLAN.md}) are realized as
  $\frac s4\to\frac s8\to\frac{3s}{32}\to\frac s{16}\to\frac s{32}$ with $s=\min(s_0,1)$.
\item The provider of G3 depends only on Section~\ref{bp:p4:sec-g3} and can start at once.
\end{enumerate}

\paragraph{Ledger candidates.} Unverified until checked by an independent verifier; quotations are
from the pinned arXiv PDF of Pos01.
\begin{enumerate}[label=(L\arabic*)]
\item E-Pos01-3 (existing candidate, checked by this writer against the pinned PDF). Pos01, p.~13:
  ``the measure of the complement of all those tori in the phase space is bounded by a constant
  times the measure of $\Omega_\alpha\times T^n$, hence is $O(\alpha)$''. It should be
  $(\Omega-\Omega_\alpha)\times T^n$, as in Proposition~2. Typo.
\item New candidate. Pos01, p.~6: ``$m(\Omega-\Omega_\alpha)=O(\alpha)$, provided the boundary of
  $\Omega$ is piecewise smooth, or at least of dimension $n-1$ so that the measure of a boundary
  layer of size $\alpha$ is $O(\alpha)$''. A boundary of (Hausdorff) dimension $n-1$ does not give
  $O(\alpha)$ layers: for $n=1$ and $\Omega=(0,1)\setminus\{1/k:k\ge2\}$, $\partial\Omega$ is countable
  (dimension $0=n-1$), but the inner layer of width $\alpha$ has measure at least $c\sqrt\alpha$
  (about $(2\alpha)^{-1/2}$ points of $\partial\Omega$ are more than $2\alpha$ apart). Correction: require
  finite upper $(n-1)$-dimensional Minkowski content of $\partial\Omega$, or directly the layer bound,
  as HEAD does. Kind: imprecise hypothesis; no effect on our statements.
\item New candidate. Pos01, p.~11, proof of Proposition~2:
  ``$m(\bar\varphi(\Omega-\Omega_\alpha))\le\abs{\bar\varphi}_{L,\Omega}\,m(\Omega-\Omega_\alpha)$''. In
  dimension $n$ the Lipschitz constant enters to the power $n$ (cf.\ BC17, footnote~32, p.~18), so the
  bound should read $\abs{\bar\varphi}^n_{L,\Omega}$. Typo; only constants change.
\end{enumerate}

\paragraph{Deviations introduced here.}
\begin{enumerate}[label=(V\arabic*)]
\item Measure step in action coordinates; this refines D-MEAS-1. The displacement $G_0$ of
  Definition~\ref{bp:p4:good-actions}, which is $k_{\nabla h(I)}(\theta)-(\theta,I)$ on
  $\R^n\times A_\gamma$ and zero on $\R^n\times Z$, extended by a clamped McShane extension, replaces
  Pöschel's extension of
  $\varphi-\mathrm{id}$ in frequency space; the boundary layer of width $\gamma$ is taken in $B$, where
  HEAD's hypothesis lives, so $\nabla h(B)$ needs no regular boundary
  (Remark~\ref{bp:p4:omega-gamma}).
\item Frequencies need no distance to the boundary of $\nabla h(B)$: $\Omega_\gamma=\nabla h(B)\cap\Dio_{\gamma,\tau}$
  (Remark~\ref{bp:p4:omega-gamma}).
\item Lipschitz step (refines D-LIP-1): stability is always measured from an unperturbed torus
  (Theorem~\ref{bp:p4:stab}); the phase given by G2 is bounded through the normalization by
  $\sup\abs{\tilde K-K'}$ (Step~\ref{step:p4:lip-phase}); far pairs use the a-priori bound with the
  threshold $\lambda_*\gamma^2$.
\item The Lagrangian property is proved directly in coordinates (Lemma~\ref{bp:p4:lagrangian})
  rather than through a primitive of the symplectic form as in GJLV05 Lemma~19.
\end{enumerate}

\paragraph{The frozen statements.} No defect was found in the frozen statements of HEAD, HEAD-Q
and G3: the proofs above establish them exactly as frozen. Two remarks for the record: the constants
$C,\varepsilon_0$ of HEAD depend on the witnesses $L$ and $C_B$, which the existential hypotheses
provide before $C,\varepsilon_0$ are chosen; and G3 allows a negative constant, which is harmless
(Lemma~\ref{bp:p4:preimage-measure} uses $\max(C_3,0)$).
